\chapter{Polynomial Response Profiles for Relative Mealy Actions}
\label{ch:polynomial-relative-mealy-actions}

\section{Overview}
\label{sec:poly-relative-overview}

This chapter develops the polynomial, or total-response, presentation of the
relative Mealy semantics constructed in
Chapter~\ref{ch:relative-mealy-program-categories}.  The span semantics there
is the uncurried execution semantics of a relative Eilenberg--Moore action.
The polynomial semantics here is its dependent-product transpose.  The use of
polynomials in the form
\[
I\xleftarrow{s}E\xrightarrow{p}B\xrightarrow{t}J
\]
and of the associated functor
\[
\Sigma_t\Pi_p s^\ast
\]
follows the standard dependent polynomial-functor calculus over locally
cartesian closed categories \cite{GambinoKock2013Polynomial}.  The language of
policies, positions, and total response profiles is also close to the
shape-and-position reading of containers
\cite{AbbottAltenkirchGhani2005Containers}.

Let
\[
\Phi=(P,\delta_P,\omega_P):\mathbb A\rightsquigarrow\mathbb B
\]
be a Mealy morphism with carrier span
\[
\begin{tikzcd}
& P \arrow[dl,"p_P"'] \arrow[dr,"q_P"] & \\
A_0 && B_0 .
\end{tikzcd}
\]
Let
\[
Y=(Y\xrightarrow{r}B_0,\sigma)
\]
be a right \(\mathbb B\)-action.  The relative state object is the pullback
\[
\begin{tikzcd}
S_{\Phi,Y}
  \arrow[r,"\pi_Y"]
  \arrow[d,"\pi_P"']
&
Y
  \arrow[d,"r"]
\\
P
  \arrow[r,"q_P"']
&
B_0
\arrow[ul,phantom,very near start,"\lrcorner"]
\end{tikzcd}
\]
so that
\[
S_{\Phi,Y}=P\times_{q_P,B_0,r}Y.
\]
The restricted right \(\mathbb A\)-action
\[
\Phi^\ast Y
\]
has action map
\[
\alpha_{\Phi,Y}:
A_1\times_{t_A,A_0,p_P\pi_P}S_{\Phi,Y}
\to S_{\Phi,Y}
\]
given by
\[
\alpha_{\Phi,Y}(a,x,y)
=
\bigl(
\delta_P(a,x),
\sigma(\omega_P(a,x),y)
\bigr).
\]
The associated action category is the relative program category
\[
\mathsf{Prog}_{\Phi,Y}
=
\int_{\mathbb A}\Phi^\ast Y.
\]

The previous chapter also constructed the downstream comparison functor
\[
\Lambda_{\Phi,Y}:
\mathsf{Prog}_{\Phi,Y}\to\int_{\mathbb B}Y,
\]
whose arrow map records the emitted \(\mathbb B\)-arrow as the downstream
transition
\[
(a,x,y)\longmapsto(\omega_P(a,x),y).
\]
Thus the labelled one-step execution data is
\[
(a,x,y)
\longmapsto
\left(
\omega_P(a,x),
\left(
\delta_P(a,x),
\sigma(\omega_P(a,x),y)
\right)
\right).
\]

Whenever the relevant dependent products exist, this uncurried labelled action
transposes to a total response profile:
\[
(x,y)
\longmapsto
\left(
a
\longmapsto
\left(
\omega_P(a,x),
\left(
\delta_P(a,x),
\sigma(\omega_P(a,x),y)
\right)
\right)
\right).
\]
The polynomial coalgebra is this dependent-product transpose.

The guiding slogan is
\[
\boxed{
\text{span semantics is uncurried execution semantics; polynomial semantics is
curried total-response semantics.}
}
\]

The second point of the chapter is variance.  A Mealy simulation
\[
\Theta:\Phi\Rightarrow\Psi
\]
is not a strict output-preserving map.  It is represented by an
output-comparison Kleisli cell
\[
Q\to B\circ P.
\]
In components, for
\[
\Phi=(P,\delta_P,\omega_P),
\qquad
\Psi=(Q,\delta_Q,\omega_Q),
\]
it has the form
\[
\theta:Q\to P,
\qquad
\alpha:Q\to B_1,
\]
with
\[
p_P\theta=p_Q,
\qquad
s_B\alpha=q_P\theta,
\qquad
t_B\alpha=q_Q.
\]
The component data are displayed by
\[
\begin{tikzcd}[column sep=large,row sep=large]
&&
B_1
  \arrow[dd,bend left=18,"s_B"]
  \arrow[dd,bend right=18,"t_B"']
&
\\
Q
  \arrow[rru,"\alpha"]
  \arrow[r,"\theta"']
  \arrow[drr,"q_Q"']
  \arrow[ddrr,bend right=16,"p_Q"']
&
P
  \arrow[dr,"q_P"]
  \arrow[rdd,bend left=16,"p_P"]
&&
\\
&&
B_0
&
\\
&&
A_0
&
\end{tikzcd}
\]
and each \(\alpha_y\) is an arrow
\[
q_P(\theta y)\xrightarrow{\alpha_y}q_Q(y)
\]
of \(\mathbb B\).

After evaluation in a downstream action \(Y\), the simulation gives
\[
h_{\Theta,Y}:S_{\Psi,Y}\to S_{\Phi,Y},
\qquad
h_{\Theta,Y}(y,z)
=
\bigl(\theta y,\sigma(\alpha_y,z)\bigr).
\]
This map is strict over \(A_0\), but not strict over the polynomial response
interface
\[
J_Y=A_0\times Y.
\]
Therefore Mealy simulations do not become ordinary coalgebra morphisms over
\(J_Y\).  They become input-strict, output-lax profile comparisons.  This
distinction is analogous in spirit to coalgebraic settings where coalgebras
live in Kleisli or Eilenberg--Moore contexts rather than as ordinary
coalgebras over a base category
\cite{BeoharKonigKupperMikaMichalski2022LiftingsSideEffects}.

The second guiding slogan is
\[
\boxed{
\text{Mealy simulations evaluate to output-lax comparisons of polynomial
response profiles.}
}
\]

Flattening turns these output-lax profile comparisons back into the
input-strict, downstream-lax labelled functors
\[
H_{\Theta,Y}:
\mathsf{Prog}_{\Psi,Y}\to\mathsf{Prog}_{\Phi,Y}
\]
constructed in Chapter~\ref{ch:relative-mealy-program-categories}.

\section{Standing hypotheses and categorical levels}
\label{sec:poly-relative-hypotheses}

\subsection{Levels of structure}
\label{subsec:poly-relative-levels}

The constructions of this chapter occur at several levels of categorical
strength:
\[
\begin{array}{c|c|c}
\text{level} & \text{categorical structure} & \text{construction} \\
\hline
0 & \text{finite limits} &
\Phi^\ast Y,\ \mathsf{Prog}_{\Phi,Y},\ J_Y=A_0\times Y \\
1 & \text{relevant dependent products} &
\text{polynomial transposes of labelled actions} \\
2 & \text{regular images} &
\diamond\text{-modalities and existential profile liftings} \\
3 & \text{coherent finite joins} &
\text{finite guarded choice} \\
4 & \text{right adjoints to pullback} &
\Box\text{-modalities and profile boxes}.
\end{array}
\]

The relative Eilenberg--Moore and program-category constructions require
finite limits.  The polynomial layer is additional and requires the specific
dependent products appearing below.  A sufficient global hypothesis is that
\(\mathcal E\) is locally cartesian closed, hence in particular finitely
complete, and that the dependent products used below exist.  This is the
standard setting for the polynomial-functor calculus
\cite{GambinoKock2013Polynomial}.

For the modal material in this chapter, we use the regular, coherent, and
first-order Heyting structure needed for diamonds, finite joins, and boxes.

\subsection{Dependent sums, pullbacks, and products}
\label{subsec:poly-relative-dependent-operators}

For a morphism
\[
f:X\to Y,
\]
write
\[
f^\ast:\mathcal E/Y\to\mathcal E/X
\]
for pullback along \(f\),
\[
\Sigma_f:\mathcal E/X\to\mathcal E/Y
\]
for postcomposition with \(f\), and
\[
\Pi_f:\mathcal E/X\to\mathcal E/Y
\]
for the right adjoint to \(f^\ast\), when it exists:
\[
\Sigma_f\dashv f^\ast\dashv\Pi_f.
\]
The adjunctions are summarized by
\[
\begin{tikzcd}[column sep=large]
\mathcal E/X
  \arrow[r,shift left=1.1ex,"\Sigma_f"]
&
\mathcal E/Y
  \arrow[l,shift left=.2ex,"f^\ast"]
  \arrow[l,shift right=1.1ex,"\Pi_f"']
\end{tikzcd}
\]
where
\[
\Sigma_f\dashv f^\ast
\qquad\text{and}\qquad
f^\ast\dashv\Pi_f.
\]

All equalities of spans and predicate transformers are read after the
canonical associativity, unit, distributivity, and Beck--Chevalley
identifications.

\section{Polynomial spans and linear spans}
\label{sec:poly-relative-polynomial-spans}

\subsection{Polynomial spans}
\label{subsec:poly-relative-polynomial-spans}

\begin{definition}[Polynomial span]
\label{def:poly-relative-polynomial-span}
A \textbf{polynomial span} from \(I\) to \(J\) is a diagram
\[
\begin{tikzcd}[column sep=large]
I
&
E
  \arrow[l,"s"']
  \arrow[r,"p"]
&
B
  \arrow[r,"t"]
&
J.
\end{tikzcd}
\]
When the required dependent product exists, it induces a functor
\[
\Sigma_t\Pi_p s^\ast:
\mathcal E/I\to\mathcal E/J.
\]
\end{definition}

Thus a polynomial span is read as the composite
\[
\begin{tikzcd}[column sep=huge]
\mathcal E/I
  \arrow[r,"s^\ast"]
&
\mathcal E/E
  \arrow[r,"\Pi_p"]
&
\mathcal E/B
  \arrow[r,"\Sigma_t"]
&
\mathcal E/J .
\end{tikzcd}
\]
This is the dependent polynomial functor associated to the displayed diagram
\cite{GambinoKock2013Polynomial}.

Fibrewise, if \(X\to I\), then
\[
\left(\Sigma_t\Pi_p s^\ast X\right)_j
\cong
\sum_{b\in B_j}
\prod_{e\in E_b}
X_{s(e)}.
\]

\subsection{Linear spans}
\label{subsec:poly-relative-linear-spans}

\begin{remark}[Linear spans]
\label{rem:poly-relative-linear-spans}
An ordinary span
\[
I\xleftarrow{s}R\xrightarrow{t}J
\]
is the linear polynomial span
\[
\begin{tikzcd}[column sep=large]
I
&
R
  \arrow[l,"s"']
  \arrow[r,"1_R"]
&
R
  \arrow[r,"t"]
&
J.
\end{tikzcd}
\]
The induced functor is
\[
\Sigma_t\Pi_{1_R}s^\ast
\cong
\Sigma_t s^\ast.
\]
Thus ordinary spans are the special case in which each execution witness has a
single continuation position.  The response polynomials below replace this
single position by a dependent family of admissible input-query positions.
\end{remark}

The polynomial-functor reading of a span is covariant:
\[
\Sigma_t s^\ast:\mathcal E/I\to\mathcal E/J.
\]
The dynamic-logic reading used later is the predicate-transformer reading:
\[
\langle R\rangle=\exists_s t^\ast.
\]
The two readings use the same span shape but act on different indexed
categories and with different variance.

\section{Currying right internal actions}
\label{sec:poly-relative-currying-actions}

\subsection{Right actions as uncurried execution maps}
\label{subsec:poly-relative-actions-uncurried}

Let
\[
\mathbb A=(A_0,A_1,s_A,t_A,e_A,m_A)
\]
be an internal category in \(\mathcal E\).  A right \(\mathbb A\)-action
\[
(X\xrightarrow{p}A_0,\rho)
\]
has action map
\[
\rho:
A_1\times_{t_A,A_0,p}X\to X
\]
satisfying
\[
p\rho=s_A\pi_1.
\]
The typing is displayed by
\[
\begin{tikzcd}[column sep=large,row sep=large]
A_1\times_{t_A,A_0,p}X
  \arrow[rr,"\rho"]
  \arrow[dr,"s_A\pi_1"']
&&
X
  \arrow[dl,"p"]
\\
&
A_0
&
\end{tikzcd}
\]
and the action domain is the pullback
\[
\begin{tikzcd}
A_1\times_{t_A,A_0,p}X
  \arrow[r,"\pi_X"]
  \arrow[d,"\pi_1"']
&
X
  \arrow[d,"p"]
\\
A_1
  \arrow[r,"t_A"']
&
A_0
\arrow[ul,phantom,very near start,"\lrcorner"] .
\end{tikzcd}
\]

Equivalently, writing
\[
t_A^\ast X=A_1\times_{t_A,A_0,p}X,
\qquad
s_A^\ast X=A_1\times_{s_A,A_0,p}X,
\]
the action map is a morphism over \(A_1\)
\[
t_A^\ast X\to s_A^\ast X.
\]
The two pullbacks sit over \(A_1\) as
\[
\begin{tikzcd}[column sep=large,row sep=large]
t_A^\ast X
  \arrow[rr,dashed,"{(\pi_1,\rho)}"]
  \arrow[dr,"\pi_X"']
  \arrow[ddr,bend right=14,"\pi_1"']
&&
s_A^\ast X
  \arrow[dl,"\pi_X"]
  \arrow[ddl,bend left=14,"\pi_1"]
\\
&
X
  \arrow[d,"p"]
&
\\
&
A_1
&
\end{tikzcd}
\]
where the dashed arrow is the action map together with the identity on the
\(A_1\)-component.

\subsection{The action-response polynomial}
\label{subsec:poly-relative-action-response-polynomial}

Assume that the dependent product
\[
\Pi_{t_A}
\]
exists.

\begin{definition}[Action-response polynomial]
\label{def:poly-relative-action-response-polynomial}
The \textbf{action-response polynomial} of \(\mathbb A\) is the endofunctor on
\(\mathcal E/A_0\)
\[
\mathcal R_{\mathbb A}(X)
:=
\Pi_{t_A}s_A^\ast X.
\]
Fibrewise,
\[
\mathcal R_{\mathbb A}(X)_v
\cong
\prod_{a:u\to v}X_u.
\]
\end{definition}

The action map transposes along
\[
t_A^\ast\dashv\Pi_{t_A}
\]
to a map over \(A_0\)
\[
\widehat\rho:
X\to\mathcal R_{\mathbb A}(X).
\]
The transpose relation is displayed by
\[
\begin{tikzcd}[column sep=large,row sep=large]
t_A^\ast X
  \arrow[rr,"\rho"]
  \arrow[d]
&&
s_A^\ast X
  \arrow[d]
\\
A_1
  \arrow[rr,equal]
&&
A_1
\end{tikzcd}
\qquad\Longleftrightarrow\qquad
\begin{tikzcd}[column sep=large,row sep=large]
X
  \arrow[rr,"\widehat\rho"]
  \arrow[dr,"p"']
&&
\Pi_{t_A}s_A^\ast X
  \arrow[dl]
\\
&
A_0
&
\end{tikzcd}
\]
and in generalized-element notation
\[
\widehat\rho(x)(a)=\rho(a,x).
\]

\subsection{Multiplicativity}
\label{subsec:poly-relative-action-response-multiplicativity}

\begin{definition}[Multiplicative action-response coalgebra]
\label{def:poly-relative-multiplicative-action-coalgebra}
A coalgebra
\[
c:X\to\mathcal R_{\mathbb A}(X)
\]
is \textbf{multiplicative} if, writing
\[
c(x)(a)=x_a,
\]
it satisfies
\[
c(x)(e_A(px))=x,
\]
and, for
\[
a:u\to v,
\qquad
b:v\to w,
\qquad
p(x)=w,
\]
if
\[
c(x)(b)=x_b,
\qquad
c(x_b)(a)=x_{ab},
\]
then
\[
c(x)(m_A(a,b))=x_{ab}.
\]
\end{definition}

Here
\[
m_A(a,b)=b\circ a.
\]
Operationally, the right action executes \(b\) first at the current state and
then executes \(a\) at the updated state.

\begin{proposition}[Actions as multiplicative response coalgebras]
\label{prop:poly-relative-actions-as-response-coalgebras}
Assume \(\Pi_{t_A}\) exists.  To give a right \(\mathbb A\)-action
\[
(X\xrightarrow{p}A_0,\rho)
\]
is equivalently to give a multiplicative coalgebra
\[
c:X\to\mathcal R_{\mathbb A}(X).
\]
\end{proposition}

\begin{proof}
\leavevmode
\begin{enumerate}
\item \textbf{Transpose the action map.}
The adjunction
\[
t_A^\ast\dashv\Pi_{t_A}
\]
identifies maps
\[
t_A^\ast X\to s_A^\ast X
\]
over \(A_1\) with maps
\[
X\to\Pi_{t_A}s_A^\ast X
\]
over \(A_0\).

\item \textbf{Read the transpose in elements.}
Under this identification,
\[
\rho(a,x)
\]
corresponds to
\[
c(x)(a).
\]

\item \textbf{Compare unit laws.}
The action unit law
\[
\rho(e_A(px),x)=x
\]
is exactly
\[
c(x)(e_A(px))=x.
\]

\item \textbf{Compare multiplication laws.}
For
\[
a:u\to v,
\qquad
b:v\to w,
\qquad
p(x)=w,
\]
the action multiplication law
\[
\rho(m_A(a,b),x)=\rho(a,\rho(b,x))
\]
is exactly the multiplicativity equation
\[
c(x)(m_A(a,b))=c(c(x)(b))(a).
\]
\end{enumerate}
\end{proof}

\section{Labelled relative actions}
\label{sec:poly-relative-labelled-actions}

\subsection{Labelled actions over a downstream action}
\label{subsec:poly-relative-labelled-actions-definition}

Fix an internal category \(\mathbb B\) and a right \(\mathbb B\)-action
\[
Y=(Y\xrightarrow{r}B_0,\sigma).
\]
The action category of \(Y\) has arrow object
\[
B_1\times_{t_B,B_0,r}Y
\]
and source-target span
\[
\begin{tikzcd}
&
B_1\times_{t_B,B_0,r}Y
  \arrow[dl,"\pi_Y"']
  \arrow[dr,"\sigma"]
&
\\
Y && Y .
\end{tikzcd}
\]

\begin{definition}[Labelled relative action]
\label{def:poly-relative-labelled-action}
A \textbf{labelled relative action over \(Y\)} consists of:
\begin{enumerate}
\item a right \(\mathbb A\)-action
\[
(X\xrightarrow{p_X}A_0,\rho);
\]
\item a map
\[
y_X:X\to Y;
\]
\item a map
\[
\lambda:
A_1\times_{t_A,A_0,p_X}X\to B_1
\]
such that
\[
t_B\lambda(a,z)=r(y_Xz),
\]
and
\[
y_X(\rho(a,z))
=
\sigma(\lambda(a,z),y_Xz);
\]
\item the unit and multiplication laws saying that
\[
(a,z)\longmapsto(\lambda(a,z),y_Xz)
\]
is the arrow map of an internal functor
\[
\Lambda_X:\int_{\mathbb A}X\to\int_{\mathbb B}Y.
\]
\end{enumerate}
\end{definition}

The arrow map into the downstream action category is typed by the diagram
\[
\begin{tikzcd}[column sep=large,row sep=large]
A_1\times_{t_A,A_0,p_X}X
  \arrow[rr,"{(a,z)\mapsto(\lambda(a,z),y_Xz)}"]
  \arrow[dr,"y_X\pi_2"']
&&
B_1\times_{t_B,B_0,r}Y
  \arrow[dl,"\pi_Y"]
  \arrow[d,"\pi_{B_1}" description]
\\
&
Y
  \arrow[d,"r"']
&
B_1
  \arrow[dl,"t_B"]
  \arrow[dr,"s_B"]
\\
&
B_0
&&
B_0 .
\end{tikzcd}
\]
The left triangle records the source condition
\[
t_B\lambda=r y_X\pi_2.
\]
The target preservation condition is the triangle
\[
\begin{tikzcd}[column sep=large,row sep=large]
A_1\times_{t_A,A_0,p_X}X
  \arrow[rr,"\rho"]
  \arrow[dr,"{(\lambda,y_X\pi_2)}"']
&&
X
  \arrow[dl,"y_X"]
\\
&
Y
&
\end{tikzcd}
\]
where the diagonal map lands in the downstream arrow object and is followed by
the downstream target map \(\sigma\).

Thus a labelled relative action is equivalently a right \(\mathbb A\)-action
whose action category maps to the downstream action category
\[
\int_{\mathbb B}Y.
\]

\subsection{The labelled action associated to a Mealy morphism}
\label{subsec:poly-relative-mealy-labelled-action}

For the labelled action associated to a Mealy morphism
\[
\Phi=(P,\delta_P,\omega_P):\mathbb A\rightsquigarrow\mathbb B,
\]
one has
\[
X=S_{\Phi,Y}=P\times_{B_0}Y,
\]
\[
y_X=\pi_Y:S_{\Phi,Y}\to Y,
\]
\[
\rho(a,x,y)
=
\bigl(
\delta_P(a,x),
\sigma(\omega_P(a,x),y)
\bigr),
\]
and
\[
\lambda(a,x,y)=\omega_P(a,x).
\]
The construction is displayed by
\[
\begin{tikzcd}[column sep=large,row sep=large]
A_1\times_{t_A,A_0,p_P\pi_P}S_{\Phi,Y}
  \arrow[rr,"{\rho=\alpha_{\Phi,Y}}"]
  \arrow[dr,"{(\omega_P,\pi_Y)}"']
&&
S_{\Phi,Y}
  \arrow[dl,"\pi_Y"]
\\
&
Y
&
\end{tikzcd}
\]
where the diagonal map represents the downstream arrow
\[
(a,x,y)\longmapsto(\omega_P(a,x),y).
\]

\begin{remark}
\label{rem:poly-relative-labelled-not-bare-action}
The bare action \(\Phi^\ast Y\) remembers the next coupled state
\[
\bigl(
\delta_P(a,x),
\sigma(\omega_P(a,x),y)
\bigr).
\]
It need not remember the raw output arrow \(\omega_P(a,x)\), because different
\(\mathbb B\)-arrows may induce the same transition in \(Y\).  Output-sensitive
guarded modalities therefore use the labelled relative action
\[
(\Phi^\ast Y,\Lambda_{\Phi,Y}),
\]
not only the underlying right \(\mathbb A\)-action.
\end{remark}

\section{The relative response interface}
\label{sec:poly-relative-interface}

\subsection{The interface object}
\label{subsec:poly-relative-interface-object}

Fix the downstream action
\[
Y=(Y\xrightarrow{r}B_0,\sigma).
\]

\begin{definition}[Relative response interface]
\label{def:poly-relative-response-interface}
The \textbf{relative response interface} is
\[
J_Y:=A_0\times Y.
\]
Write
\[
\pi_A:J_Y\to A_0,
\qquad
\pi_Y:J_Y\to Y
\]
for its projections.
\end{definition}

The product is displayed by
\[
\begin{tikzcd}
J_Y=A_0\times Y
  \arrow[r,"\pi_Y"]
  \arrow[d,"\pi_A"']
&
Y
  \arrow[d]
\\
A_0
  \arrow[r]
&
1
\arrow[ul,phantom,very near start,"\lrcorner"] .
\end{tikzcd}
\]

For a labelled relative action \(X\), define
\[
j_X:X\to J_Y,
\qquad
j_X(z)=(p_Xz,y_Xz).
\]
The defining diagram is
\[
\begin{tikzcd}[column sep=large,row sep=large]
&
X
  \arrow[dl,"p_X"']
  \arrow[dr,"y_X"]
  \arrow[d,dashed,"j_X" description]
&
\\
A_0
&
J_Y
  \arrow[l,"\pi_A"]
  \arrow[r,"\pi_Y"']
&
Y .
\end{tikzcd}
\]
For a Mealy morphism \(\Phi\), this gives
\[
j_{\Phi,Y}:S_{\Phi,Y}\to J_Y,
\qquad
j_{\Phi,Y}(x,y)=(p_Px,y).
\]

\subsection{Interfaces and simulations}
\label{subsec:poly-relative-interface-simulations}

\begin{remark}[Interface and simulations]
\label{rem:poly-relative-interface-and-simulations}
The interface \(J_Y=A_0\times Y\) is the base for strict response profiles.  A
Mealy simulation does not generally preserve this base strictly.  If
\[
\Theta:\Phi\Rightarrow\Psi
\]
has output-comparison component \(\alpha\), then
\[
h_{\Theta,Y}(y,z)
=
(\theta y,\sigma(\alpha_y,z))
\]
preserves the \(A_0\)-component, but changes the \(Y\)-component from \(z\) to
\[
\sigma(\alpha_y,z).
\]
This is why simulations become output-lax profile comparisons rather than
ordinary coalgebra morphisms over \(J_Y\).
\end{remark}

The failure of strictness over \(J_Y\) is displayed by
\[
\begin{tikzcd}[column sep=large,row sep=large]
S_{\Psi,Y}
  \arrow[r,"h_{\Theta,Y}"]
  \arrow[d,"j_{\Psi,Y}"']
&
S_{\Phi,Y}
  \arrow[d,"j_{\Phi,Y}"]
\\
J_Y
  \arrow[r,dashed,"?"]
&
J_Y .
\end{tikzcd}
\]
The \(A_0\)-component is preserved, but the \(Y\)-component is transformed by
the evaluated comparison arrow.

\section{Queries, responses, and policies}
\label{sec:poly-relative-queries-responses}

The following construction has the usual container shape: one first forms an
object of admissible input positions, then an object of possible responses at
each position, and finally a dependent product of responses over all positions
at a given interface \cite{AbbottAltenkirchGhani2005Containers}.

\subsection{Input queries}
\label{subsec:poly-relative-input-queries}

\begin{definition}[Input-query object]
\label{def:poly-relative-input-query-object}
The \textbf{input-query object} is the pullback
\[
\mathrm{Inp}_Y:=A_1\times_{t_A,A_0,\pi_A}J_Y.
\]
Thus
\[
\begin{tikzcd}[column sep=large,row sep=large]
\mathrm{Inp}_Y
  \arrow[r,"\overline a"]
  \arrow[d,"\kappa_Y"']
&
A_1
  \arrow[d,"t_A"]
\\
J_Y
  \arrow[r,"\pi_A"']
&
A_0
\arrow[ul,phantom,very near start,"\lrcorner"]
\end{tikzcd}
\]
is a pullback square.
\end{definition}

A generalized element of \(\mathrm{Inp}_Y\) is written
\[
(a,y),
\]
where \(a:u\to v\) is an arrow of \(\mathbb A\) and \(y\in Y\).  Its current
interface is
\[
\kappa_Y(a,y)=(t_Aa,y).
\]
There is also a next-input-interface map
\[
\lambda_A:\mathrm{Inp}_Y\to A_0,
\qquad
\lambda_A(a,y)=s_A(a).
\]
The maps are displayed by
\[
\begin{tikzcd}[column sep=large,row sep=large]
&
\mathrm{Inp}_Y
  \arrow[dl,"\lambda_A"']
  \arrow[d,"\kappa_Y"]
  \arrow[dr,"\overline a"]
&
\\
A_0
&
J_Y
  \arrow[l,"\pi_A"]
  \arrow[r,"\pi_Y"']
&
A_1
  \arrow[l,bend left=12,"t_A"]
  \arrow[ll,bend right=16,"s_A"']
\end{tikzcd}
\]

\subsection{Output responses}
\label{subsec:poly-relative-output-responses}

A response to a query \((a:u\to v,y)\) is an arrow of \(\mathbb B\)
\[
\beta:b'\to r(y).
\]
This arrow may be consumed by the right \(\mathbb B\)-action \(Y\), giving the
downstream continuation state
\[
\sigma(\beta,y).
\]

\begin{definition}[Output-response object]
\label{def:poly-relative-output-response-object}
The \textbf{output-response object} is
\[
\mathrm{Resp}_Y
:=
\mathrm{Inp}_Y\times_{r\pi_Y\kappa_Y,B_0,t_B}B_1.
\]
Thus
\[
\begin{tikzcd}[column sep=large,row sep=large]
\mathrm{Resp}_Y
  \arrow[r,"\overline\beta"]
  \arrow[d,"q_Y"']
&
B_1
  \arrow[d,"t_B"]
\\
\mathrm{Inp}_Y
  \arrow[r,"r\pi_Y\kappa_Y"']
&
B_0
\arrow[ul,phantom,very near start,"\lrcorner"]
\end{tikzcd}
\]
is a pullback square.
\end{definition}

A generalized element is written
\[
(a,y,\beta),
\]
where
\[
t_B(\beta)=r(y).
\]

Define the continuation-interface map
\[
n_Y:\mathrm{Resp}_Y\to J_Y
\]
by
\[
n_Y(a,y,\beta)
=
\bigl(
s_A(a),
\sigma(\beta,y)
\bigr).
\]
Equivalently,
\[
\pi_A n_Y=\lambda_A q_Y,
\qquad
\pi_Y n_Y=\sigma(\overline\beta,\pi_Y\kappa_Yq_Y).
\]
The typing of \(n_Y\) is displayed by
\[
\begin{tikzcd}[column sep=large,row sep=large]
\mathrm{Resp}_Y
  \arrow[rr,"n_Y"]
  \arrow[drr,bend left=16,"\lambda_Aq_Y"]
  \arrow[dr,"\overline\beta"']
&&
J_Y
  \arrow[d,"\pi_Y"]
  \arrow[ddl,bend left=16,"\pi_A"]
\\
&
B_1
&
Y
\\
&
A_0
&
\end{tikzcd}
\]
where the \(Y\)-component is obtained by applying the action map \(\sigma\).

\subsection{Total response policies}
\label{subsec:poly-relative-total-policies}

A total output policy at an interface
\[
(v,y)\in J_Y
\]
assigns to every input arrow
\[
a:u\to v
\]
an output arrow
\[
\beta_a:b'_a\to r(y).
\]

\begin{definition}[Policy object]
\label{def:poly-relative-policy-object}
Assume the dependent product below exists.  The \textbf{policy object} is
\[
K_{\mathbb A,\mathbb B,Y}
:=
\Pi_{\kappa_Y}(\mathrm{Resp}_Y\xrightarrow{q_Y}\mathrm{Inp}_Y).
\]
It is an object over \(J_Y\).  Write
\[
t_{\mathbb A,\mathbb B,Y}:
K_{\mathbb A,\mathbb B,Y}\to J_Y
\]
for its structure map.
\end{definition}

A generalized element of \(K_{\mathbb A,\mathbb B,Y}\) over \((v,y)\) is a
policy
\[
\pi_0:
\{a:u\to v\}
\to
\{\beta:b'\to r(y)\}.
\]

Now form the pullback
\[
E_{\mathbb A,\mathbb B,Y}
:=
\mathrm{Inp}_Y\times_{J_Y}K_{\mathbb A,\mathbb B,Y},
\]
where the maps to \(J_Y\) are \(\kappa_Y\) and
\(t_{\mathbb A,\mathbb B,Y}\).  Thus
\[
\begin{tikzcd}[column sep=large,row sep=large]
E_{\mathbb A,\mathbb B,Y}
  \arrow[r,"\epsilon_Q"]
  \arrow[d,"p_{\mathbb A,\mathbb B,Y}"']
&
\mathrm{Inp}_Y
  \arrow[d,"\kappa_Y"]
\\
K_{\mathbb A,\mathbb B,Y}
  \arrow[r,"t_{\mathbb A,\mathbb B,Y}"']
&
J_Y
\arrow[ul,phantom,very near start,"\lrcorner"]
\end{tikzcd}
\]
is a pullback square.

The counit of
\[
\kappa_Y^\ast\dashv\Pi_{\kappa_Y}
\]
gives an evaluation map
\[
\operatorname{ev}:
E_{\mathbb A,\mathbb B,Y}\to \mathrm{Resp}_Y
\]
over \(\mathrm{Inp}_Y\).  The evaluation diagram is
\[
\begin{tikzcd}[column sep=large,row sep=large]
E_{\mathbb A,\mathbb B,Y}
  \arrow[rr,"\operatorname{ev}"]
  \arrow[dr,"\epsilon_Q"']
&&
\mathrm{Resp}_Y
  \arrow[dl,"q_Y"]
\\
&
\mathrm{Inp}_Y
&
\end{tikzcd}
\]
and in generalized-element notation,
\[
\operatorname{ev}(a,y,\pi_0)=(a,y,\pi_0(a)).
\]

Define
\[
s_{\mathbb A,\mathbb B,Y}
:=
n_Y\operatorname{ev}:
E_{\mathbb A,\mathbb B,Y}\to J_Y.
\]
The structure maps are summarized by
\[
\begin{tikzcd}[column sep=large,row sep=large]
&
E_{\mathbb A,\mathbb B,Y}
  \arrow[dl,"s_{\mathbb A,\mathbb B,Y}"']
  \arrow[d,"p_{\mathbb A,\mathbb B,Y}"]
  \arrow[dr,"\operatorname{ev}"]
&
\\
J_Y
&
K_{\mathbb A,\mathbb B,Y}
  \arrow[l,"t_{\mathbb A,\mathbb B,Y}"]
&
\mathrm{Resp}_Y
  \arrow[ll,"n_Y"]
\end{tikzcd}
\]

\subsection{Dependent products required}
\label{subsec:poly-relative-dependent-products-required}

We also assume that the projection
\[
p_{\mathbb A,\mathbb B,Y}:
E_{\mathbb A,\mathbb B,Y}\to K_{\mathbb A,\mathbb B,Y}
\]
admits a dependent product.  This is automatic if \(\mathcal E\) is locally
cartesian closed.  More generally, it follows whenever \(\kappa_Y\) is
exponentiable and exponentiable morphisms are stable under pullback, since
\(p_{\mathbb A,\mathbb B,Y}\) is the pullback of \(\kappa_Y\) along
\[
t_{\mathbb A,\mathbb B,Y}:K_{\mathbb A,\mathbb B,Y}\to J_Y.
\]

\section{The relative Mealy response polynomial}
\label{sec:poly-relative-response-polynomial}

\subsection{Definition of the polynomial}
\label{subsec:poly-relative-response-polynomial-definition}

\begin{definition}[Relative Mealy response polynomial]
\label{def:poly-relative-mealy-response-polynomial}
The \textbf{relative Mealy response polynomial} associated to
\[
(\mathbb A,\mathbb B,Y)
\]
is the polynomial span
\[
\begin{tikzcd}[column sep=large]
J_Y
&
E_{\mathbb A,\mathbb B,Y}
  \arrow[l,"s_{\mathbb A,\mathbb B,Y}"']
  \arrow[r,"p_{\mathbb A,\mathbb B,Y}"]
&
K_{\mathbb A,\mathbb B,Y}
  \arrow[r,"t_{\mathbb A,\mathbb B,Y}"]
&
J_Y.
\end{tikzcd}
\]
It induces an endofunctor
\[
\mathcal M_{\mathbb A,\mathbb B,Y}:
\mathcal E/J_Y\to\mathcal E/J_Y
\]
defined by
\[
\mathcal M_{\mathbb A,\mathbb B,Y}(X)
=
\Sigma_{t_{\mathbb A,\mathbb B,Y}}
\Pi_{p_{\mathbb A,\mathbb B,Y}}
s_{\mathbb A,\mathbb B,Y}^{\ast}X.
\]
\end{definition}

In the following diagram,
\[
J_Y^{\mathrm{cur}}
\qquad\text{and}\qquad
J_Y^{\mathrm{next}}
\]
are two displayed copies of the same object \(J_Y\), used only to distinguish
the current response interface from the continuation response interface.

The full construction can be displayed as
\[
\begin{tikzcd}[column sep=large,row sep=large]
&
\mathrm{Resp}_Y
  \arrow[rr,"\overline\beta"]
  \arrow[d,"q_Y"']
  \arrow[ddd,bend left=18,"n_Y"]
&&
B_1
  \arrow[d,"t_B"]
\\
E_{\mathbb A,\mathbb B,Y}
  \arrow[ur,"\operatorname{ev}"]
  \arrow[r,"\epsilon_Q" description]
  \arrow[d,"p_{\mathbb A,\mathbb B,Y}"']
  \arrow[ddr,bend right=16,"s_{\mathbb A,\mathbb B,Y}=n_Y\operatorname{ev}"']
  \arrow[dr,phantom,very near start,"\lrcorner"]
&
\mathrm{Inp}_Y
  \arrow[rr,"r\pi_Y\kappa_Y" description]
  \arrow[d,"\kappa_Y"]
  \arrow[dr,"\overline a" description]
&&
B_0
\\
K_{\mathbb A,\mathbb B,Y}
  \arrow[r,"t_{\mathbb A,\mathbb B,Y}"']
&
J_Y^{\mathrm{cur}}
&
A_1
&
\\
&
J_Y^{\mathrm{next}}
&
&
\end{tikzcd}
\]

\subsection{Dependent distributivity}
\label{subsec:poly-relative-dependent-distributivity}

\begin{lemma}[Dependent distributivity for response policies]
\label{lem:poly-relative-dependent-distributivity}
In the locally cartesian closed setting, the dependent distributivity
isomorphism for the polynomial data gives a canonical fibrewise isomorphism
\[
\sum_{\pi_0\in\prod_a O_a}
\prod_a X_{a,\pi_0(a)}
\cong
\prod_a
\sum_{\beta\in O_a}X_{a,\beta}.
\]
Equivalently, this is the internal distributivity of dependent sums over
dependent products associated to the family of response objects over the input
query object.  It is not an external choice principle.
\end{lemma}

\begin{proof}
\leavevmode
\begin{enumerate}
\item \textbf{Define the forward map.}
Given a pair
\[
(\pi_0,(x_a)_a)
\]
where
\[
\pi_0(a)\in O_a
\qquad\text{and}\qquad
x_a\in X_{a,\pi_0(a)},
\]
send it to the family
\[
\bigl(\pi_0(a),x_a\bigr)_a
\in
\prod_a\sum_{\beta\in O_a}X_{a,\beta}.
\]

\item \textbf{Define the inverse map.}
Given a family
\[
(\beta_a,x_a)_a
\]
with
\[
\beta_a\in O_a
\qquad\text{and}\qquad
x_a\in X_{a,\beta_a},
\]
define
\[
\pi_0(a):=\beta_a.
\]
Then the second components form a family
\[
x_a\in X_{a,\pi_0(a)}.
\]

\item \textbf{Check the two composites.}
The two composites are the identity by the uniqueness clauses for dependent
sums and dependent products.
\end{enumerate}
\end{proof}

\subsection{Fibre calculation}
\label{subsec:poly-relative-fibre-calculation}

\begin{proposition}[Fibre calculation]
\label{prop:poly-relative-fibre-calculation}
Let
\[
X\to J_Y=A_0\times Y
\]
be an object over the response interface.  For every generalized interface
\[
(v,y)\in J_Y,
\]
there is a canonical fibrewise description
\[
\mathcal M_{\mathbb A,\mathbb B,Y}(X)_{(v,y)}
\cong
\prod_{a:u\to v}
\sum_{\beta:b'\to r(y)}
X_{(u,\sigma(\beta,y))}.
\]
\end{proposition}

\begin{proof}
\leavevmode
\begin{enumerate}
\item \textbf{Describe policies.}
A point of \(K_{\mathbb A,\mathbb B,Y}\) over \((v,y)\) is a total policy
\[
\pi_0:a\mapsto\beta_a
\]
assigning to each input arrow \(a:u\to v\) an output arrow
\[
\beta_a:b'_a\to r(y).
\]

\item \textbf{Describe positions over a policy.}
The object \(E_{\mathbb A,\mathbb B,Y}\) has, over such a policy, one position
for each input arrow
\[
a:u\to v.
\]
The continuation-interface map sends that position to
\[
(u,\sigma(\beta_a,y)).
\]

\item \textbf{Apply the dependent product.}
Thus
\[
\Pi_{p_{\mathbb A,\mathbb B,Y}}
s_{\mathbb A,\mathbb B,Y}^{\ast}X
\]
assigns to a policy \(\pi_0\) a family
\[
x_a\in X_{(u,\sigma(\pi_0(a),y))}
\]
for every input arrow \(a:u\to v\).

\item \textbf{Apply the dependent sum.}
Applying
\[
\Sigma_{t_{\mathbb A,\mathbb B,Y}}
\]
regards the result as lying over \((v,y)\).  Therefore the fibre is first
\[
\sum_{\pi_0}
\prod_{a:u\to v}
X_{(u,\sigma(\pi_0(a),y))}.
\]

\item \textbf{Use dependent distributivity.}
The dependent distributivity isomorphism of
Lemma~\ref{lem:poly-relative-dependent-distributivity} gives
\[
\sum_{\pi_0}
\prod_{a:u\to v}
X_{(u,\sigma(\pi_0(a),y))}
\cong
\prod_{a:u\to v}
\sum_{\beta:b'\to r(y)}
X_{(u,\sigma(\beta,y))}.
\]
\end{enumerate}
\end{proof}

\begin{remark}
\label{rem:poly-relative-response-polynomial-meaning}
At interface \((v,y)\), a total response profile must respond to every
admissible input arrow \(a:u\to v\).  For each such input, it chooses an output
arrow
\[
\beta:b'\to r(y)
\]
and a continuation state over
\[
(u,\sigma(\beta,y)).
\]
\end{remark}

\section{Coalgebras and labelled one-step data}
\label{sec:poly-relative-coalgebras}

\subsection{Coalgebras over the response interface}
\label{subsec:poly-relative-coalgebras-over-interface}

Let
\[
X\xrightarrow{j_X}J_Y
\]
be an object over the response interface.  Write
\[
p_X:=\pi_Aj_X:X\to A_0,
\qquad
y_X:=\pi_Yj_X:X\to Y.
\]
Thus
\[
\begin{tikzcd}[column sep=large,row sep=large]
&
X
  \arrow[dl,"p_X"']
  \arrow[dr,"y_X"]
  \arrow[d,"j_X" description]
&
\\
A_0
&
J_Y
  \arrow[l,"\pi_A"]
  \arrow[r,"\pi_Y"']
&
Y .
\end{tikzcd}
\]

A coalgebra
\[
c:X\to\mathcal M_{\mathbb A,\mathbb B,Y}(X)
\]
assigns to each
\[
z\in X_{(v,y)}
\]
and each input arrow
\[
a:u\to v
\]
a pair
\[
c(z)(a)=(\beta_a,z_a),
\]
where
\[
\beta_a:b'_a\to r(y),
\qquad
z_a\in X_{(u,\sigma(\beta_a,y))}.
\]

Thus \(c\) determines uncurried maps
\[
\rho_c:
A_1\times_{t_A,A_0,p_X}X\to X
\]
and
\[
\lambda_c:
A_1\times_{t_A,A_0,p_X}X\to B_1
\]
by
\[
c(z)(a)=(\lambda_c(a,z),\rho_c(a,z)).
\]

These satisfy
\[
p_X\rho_c(a,z)=s_A(a),
\]
\[
t_B\lambda_c(a,z)=r(y_Xz),
\]
and
\[
y_X\rho_c(a,z)
=
\sigma(\lambda_c(a,z),y_Xz).
\]
Conversely, such maps \((\rho,\lambda)\) determine a coalgebra by
\[
c(z)(a)=(\lambda(a,z),\rho(a,z)).
\]

\subsection{Currying theorem}
\label{subsec:poly-relative-currying-labelled-one-step}

\begin{theorem}[Currying labelled one-step data]
\label{thm:poly-relative-currying-labelled-one-step}
To give a coalgebra
\[
c:X\to\mathcal M_{\mathbb A,\mathbb B,Y}(X)
\]
is equivalently to give maps
\[
\rho:
A_1\times_{t_A,A_0,p_X}X\to X,
\qquad
\lambda:
A_1\times_{t_A,A_0,p_X}X\to B_1
\]
satisfying
\[
p_X\rho=s_A\pi_1,
\]
\[
t_B\lambda=r y_X\pi_2,
\]
and
\[
y_X\rho=\sigma(\lambda,y_X\pi_2).
\]
\end{theorem}

\begin{proof}
\leavevmode
\begin{enumerate}
\item \textbf{Start from the policy fibre.}
The polynomial was defined in policy form, so its fibre over \((v,y)\) is
initially
\[
\sum_{\pi_0\in\prod_{a:u\to v}O_a}
\prod_{a:u\to v}
X_{(u,\sigma(\pi_0(a),y))},
\]
where \(O_a\) is the fibre of possible output arrows
\[
\beta:b'\to r(y)
\]
for the input \(a:u\to v\).

\item \textbf{Distribute policies into pointwise responses.}
By Lemma~\ref{lem:poly-relative-dependent-distributivity}, this is
canonically isomorphic to
\[
\prod_{a:u\to v}
\sum_{\beta:b'\to r(y)}
X_{(u,\sigma(\beta,y))}.
\]

\item \textbf{Read a coalgebra in fibres.}
Thus a coalgebra over \(J_Y\) is equivalently a choice, for every state
\[
z\in X_{(v,y)}
\]
and every admissible input
\[
a:u\to v,
\]
of an output arrow
\[
\lambda(a,z):b'\to r(y)
\]
and a continuation state
\[
\rho(a,z)\in X_{(u,\sigma(\lambda(a,z),y))}.
\]

\item \textbf{Translate the fibre conditions into morphism equations.}
The statement that
\[
\rho(a,z)\in X_{(u,\sigma(\lambda(a,z),y))}
\]
is exactly the conjunction
\[
p_X\rho=s_A\pi_1,
\]
\[
t_B\lambda=r y_X\pi_2,
\]
and
\[
y_X\rho=\sigma(\lambda,y_X\pi_2).
\]
\end{enumerate}
\end{proof}

\section{Multiplicative labelled response coalgebras}
\label{sec:poly-relative-multiplicative-coalgebras}

\subsection{Definition of multiplicativity}
\label{subsec:poly-relative-multiplicative-definition}

The coalgebra map records typed one-step labelled response data.  To correspond
to an action category, this profile must preserve identities and composition.

We keep the composition convention of the previous chapters: for internally
composable arrows
\[
x\xrightarrow{\alpha}y\xrightarrow{\beta}z
\]
in \(\mathbb B\),
\[
m_B(\alpha,\beta)=\beta\circ\alpha.
\]

\begin{definition}[Multiplicative labelled response coalgebra]
\label{def:poly-relative-multiplicative-response-coalgebra}
A coalgebra
\[
c:X\to\mathcal M_{\mathbb A,\mathbb B,Y}(X)
\]
is \textbf{multiplicative} if, writing
\[
c(z)(a)=(\beta_a,z_a),
\]
the following equations hold.
\begin{enumerate}
\item For every \(z\in X\),
\[
c(z)(e_A(p_Xz))
=
(e_B(r(y_Xz)),z).
\]

\item For every composable pair
\[
a:u\to v,
\qquad
b:v\to w,
\]
and every \(z\in X_{(w,y)}\), if
\[
c(z)(b)=(\beta,z')
\]
and
\[
c(z')(a)=(\alpha,z'')
\]
then
\[
c(z)(m_A(a,b))
=
(m_B(\alpha,\beta),z'').
\]
\end{enumerate}
\end{definition}

Here \(m_A(a,b)=b\circ a\), but operationally the right action executes \(b\)
first at the current state and then executes \(a\) at the updated state.
Similarly, the output arrow for the composite is
\[
m_B(\alpha,\beta)=\beta\circ\alpha.
\]

\subsection{Multiplicative coalgebras as labelled actions}
\label{subsec:poly-relative-coalgebras-labelled-actions}

\begin{proposition}[Multiplicative coalgebras as labelled actions]
\label{prop:poly-relative-coalgebras-labelled-actions}
Assume the relevant dependent products exist.  To give a multiplicative
coalgebra
\[
c:X\to\mathcal M_{\mathbb A,\mathbb B,Y}(X)
\]
is equivalently to give:
\begin{enumerate}
\item a right \(\mathbb A\)-action on \(X\xrightarrow{p_X}A_0\);
\item an internal functor
\[
\Lambda_X:\int_{\mathbb A}X\to\int_{\mathbb B}Y
\]
whose object map is \(y_X:X\to Y\).
\end{enumerate}
Equivalently, it is a labelled relative action over \(Y\).
\end{proposition}

\begin{proof}
\leavevmode
\begin{enumerate}
\item \textbf{Uncurry the coalgebra.}
By Theorem~\ref{thm:poly-relative-currying-labelled-one-step}, a coalgebra is
equivalent to maps
\[
\rho:
A_1\times_{t_A,A_0,p_X}X\to X,
\qquad
\lambda:
A_1\times_{t_A,A_0,p_X}X\to B_1
\]
with the displayed typing equations.

\item \textbf{Recover the right action.}
The first component \(\rho\), together with the identity and composition
clauses of multiplicativity, is exactly a right \(\mathbb A\)-action on
\(X\).

\item \textbf{Construct the action category.}
The action category \(\int_{\mathbb A}X\) has object object \(X\) and arrow
object
\[
A_1\times_{t_A,A_0,p_X}X.
\]

\item \textbf{Construct the downstream functor.}
The map
\[
(a,z)\longmapsto(\lambda(a,z),y_Xz)
\]
is well-defined as a map into the arrow object of \(\int_{\mathbb B}Y\)
because
\[
t_B\lambda(a,z)=r(y_Xz).
\]
Its source and target compatibility are displayed by
\[
\begin{tikzcd}[column sep=large,row sep=large]
A_1\times_{t_A,A_0,p_X}X
  \arrow[rr,"{(\lambda,y_X\pi_2)}"]
  \arrow[d,"\pi_2"']
&&
B_1\times_{t_B,B_0,r}Y
  \arrow[d,"\pi_Y"]
\\
X
  \arrow[rr,"y_X"']
&&
Y
\end{tikzcd}
\]
and
\[
\begin{tikzcd}[column sep=large,row sep=large]
A_1\times_{t_A,A_0,p_X}X
  \arrow[rr,"{(\lambda,y_X\pi_2)}"]
  \arrow[d,"\rho"']
&&
B_1\times_{t_B,B_0,r}Y
  \arrow[d,"\sigma"]
\\
X
  \arrow[rr,"y_X"']
&&
Y .
\end{tikzcd}
\]

\item \textbf{Compare identity and composition.}
The unit and multiplication clauses of multiplicativity are precisely identity
and composition preservation for this internal functor.
\end{enumerate}
\end{proof}

\section{The coalgebra induced by a Mealy morphism}
\label{sec:poly-relative-induced-coalgebra}

\subsection{Definition of the induced coalgebra}
\label{subsec:poly-relative-induced-coalgebra-definition}

Let
\[
\Phi=(P,\delta_P,\omega_P):\mathbb A\rightsquigarrow\mathbb B
\]
be a Mealy morphism and let
\[
Y=(Y\xrightarrow{r}B_0,\sigma)
\]
be a right \(\mathbb B\)-action.

Recall that
\[
S_{\Phi,Y}=P\times_{B_0}Y
\]
maps to
\[
J_Y=A_0\times Y
\]
by
\[
j_{\Phi,Y}(x,y)=(p_Px,y).
\]
The structure map is displayed by
\[
\begin{tikzcd}[column sep=large,row sep=large]
&
S_{\Phi,Y}
  \arrow[dl,"p_P\pi_P"']
  \arrow[dr,"\pi_Y"]
  \arrow[d,"j_{\Phi,Y}" description]
&
\\
A_0
&
J_Y
  \arrow[l,"\pi_A"]
  \arrow[r,"\pi_Y"']
&
Y .
\end{tikzcd}
\]

\begin{definition}[Relative Mealy response coalgebra]
\label{def:poly-relative-mealy-response-coalgebra}
The \textbf{relative Mealy response coalgebra} induced by \(\Phi\) at \(Y\) is
the coalgebra
\[
c_{\Phi,Y}:
S_{\Phi,Y}\to
\mathcal M_{\mathbb A,\mathbb B,Y}(S_{\Phi,Y})
\]
given fibrewise by
\[
c_{\Phi,Y}(x,y)(a)
=
\left(
\omega_P(a,x),
\left(
\delta_P(a,x),
\sigma(\omega_P(a,x),y)
\right)
\right).
\]
\end{definition}

\subsection{Transpose of the labelled action}
\label{subsec:poly-relative-mealy-coalgebra-transpose}

\begin{theorem}[Mealy coalgebra as transpose of labelled action]
\label{thm:poly-relative-mealy-coalgebra-transpose}
The coalgebra
\[
c_{\Phi,Y}
\]
is the dependent-product transpose of the labelled relative action
\[
(\Phi^\ast Y,\Lambda_{\Phi,Y}).
\]
Moreover, it is multiplicative.
\end{theorem}

\begin{proof}
\leavevmode
\begin{enumerate}
\item \textbf{Recall the action map.}
The action map of \(\Phi^\ast Y\) is
\[
(a,x,y)
\longmapsto
\left(
\delta_P(a,x),
\sigma(\omega_P(a,x),y)
\right).
\]

\item \textbf{Recall the downstream label.}
The downstream comparison functor supplies the output label
\[
(a,x,y)\longmapsto\omega_P(a,x).
\]

\item \textbf{Combine the two components.}
Together these give the uncurried labelled one-step data
\[
(a,x,y)
\longmapsto
\left(
\omega_P(a,x),
\left(
\delta_P(a,x),
\sigma(\omega_P(a,x),y)
\right)
\right).
\]

\item \textbf{Transpose.}
By Theorem~\ref{thm:poly-relative-currying-labelled-one-step}, this
uncurried labelled map transposes to \(c_{\Phi,Y}\).

\item \textbf{Check multiplicativity.}
Multiplicativity follows from the Mealy unit and multiplication equations for
\(\Phi\) and from the unit and multiplication laws of the right
\(\mathbb B\)-action \(Y\).  Equivalently, it follows from
Proposition~\ref{prop:poly-relative-coalgebras-labelled-actions}, since
\[
\Lambda_{\Phi,Y}:\mathsf{Prog}_{\Phi,Y}\to\int_{\mathbb B}Y
\]
is an internal functor.
\end{enumerate}
\end{proof}

\section{Flattening response coalgebras}
\label{sec:poly-relative-flattening}

\subsection{Position objects}
\label{subsec:poly-relative-position-objects}

Let
\[
J_Y
\xleftarrow{s}
E_{\mathbb A,\mathbb B,Y}
\xrightarrow{p}
K_{\mathbb A,\mathbb B,Y}
\xrightarrow{t}
J_Y
\]
be the relative response polynomial.

For an object
\[
X\to J_Y,
\]
define the position object at \(X\) by
\[
\mathrm{Pos}_X
:=
E_{\mathbb A,\mathbb B,Y}
\times_{K_{\mathbb A,\mathbb B,Y}}
\Pi_p(s^\ast X).
\]
It is displayed by the pullback
\[
\begin{tikzcd}[column sep=large,row sep=large]
\mathrm{Pos}_X
  \arrow[r]
  \arrow[d]
&
\Pi_p(s^\ast X)
  \arrow[d]
\\
E_{\mathbb A,\mathbb B,Y}
  \arrow[r,"p"']
&
K_{\mathbb A,\mathbb B,Y}
\arrow[ul,phantom,very near start,"\lrcorner"] .
\end{tikzcd}
\]
It carries a canonical projection
\[
\partial_X:\mathrm{Pos}_X\to
\mathcal M_{\mathbb A,\mathbb B,Y}(X).
\]
Externally, a point of \(\mathrm{Pos}_X\) is a total response profile together
with a chosen input position.  This is the usual passage from a container-like
total object to its object of positions
\cite{AbbottAltenkirchGhani2005Containers}.

The position object has evaluation maps
\[
\operatorname{in}:\mathrm{Pos}_X\to A_1,
\]
\[
\operatorname{out}:\mathrm{Pos}_X\to B_1,
\]
\[
\operatorname{down}:\mathrm{Pos}_X\to
\left(\int_{\mathbb B}Y\right)_1,
\]
and
\[
\operatorname{next}:\mathrm{Pos}_X\to X.
\]
They extract the chosen input arrow, raw output arrow, downstream transition,
and continuation state.  These maps fit into
\[
\begin{tikzcd}[column sep=large,row sep=large]
&
\mathrm{Pos}_X
  \arrow[dl,"\operatorname{in}"']
  \arrow[d,"\operatorname{down}" description]
  \arrow[dr,"\operatorname{next}"]
&
\\
A_1
&
\left(\int_{\mathbb B}Y\right)_1
  \arrow[d,"\pi_{B_1}" description]
  \arrow[r,"\sigma"]
&
Y
\\
&
B_1
&
\end{tikzcd}
\]
with
\[
\operatorname{out}
=
\pi_{B_1}\operatorname{down},
\qquad
y_X\operatorname{next}
=
\sigma\operatorname{down}.
\]

\subsection{Flattening a coalgebra}
\label{subsec:poly-relative-flattening-coalgebra}

Let
\[
c:X\to\mathcal M_{\mathbb A,\mathbb B,Y}(X)
\]
be a coalgebra.  Pulling back the position projection along \(c\) gives
\[
\begin{tikzcd}[column sep=large,row sep=large]
M_c
  \arrow[r]
  \arrow[d,"s_c"']
&
\mathrm{Pos}_X
  \arrow[d,"\partial_X"]
\\
X
  \arrow[r,"c"']
&
\mathcal M_{\mathbb A,\mathbb B,Y}(X)
\arrow[ul,phantom,very near start,"\lrcorner"] .
\end{tikzcd}
\]
The object \(M_c\) is the flattened primitive execution object of \(c\).

Under the uncurried presentation of \(c\), one has
\[
M_c\cong A_1\times_{t_A,A_0,p_X}X.
\]
The source map is
\[
s_c(a,z)=z,
\]
and the target map is
\[
t_c(a,z)=\rho_c(a,z).
\]
The flattened execution span is
\[
\begin{tikzcd}
&
M_c
  \arrow[dl,"s_c"']
  \arrow[dr,"t_c"]
&
\\
X && X .
\end{tikzcd}
\]

\subsection{Flattening multiplicative coalgebras}
\label{subsec:poly-relative-flattening-multiplicative}

\begin{theorem}[Flattening multiplicative coalgebras]
\label{thm:poly-relative-flattening-multiplicative}
If
\[
c:X\to\mathcal M_{\mathbb A,\mathbb B,Y}(X)
\]
is multiplicative, then the flattened span
\[
X\xleftarrow{s_c}M_c\xrightarrow{t_c}X
\]
is the primitive execution span of the action category
\[
\int_{\mathbb A}X.
\]
Moreover, the evaluation maps on \(M_c\) recover the input projection and the
downstream comparison functor
\[
\int_{\mathbb A}X\to\int_{\mathbb B}Y.
\]
\end{theorem}

\begin{proof}
\leavevmode
\begin{enumerate}
\item \textbf{Replace the coalgebra by labelled action data.}
By Proposition~\ref{prop:poly-relative-coalgebras-labelled-actions}, a
multiplicative coalgebra is equivalently a right \(\mathbb A\)-action on
\(X\) equipped with an internal functor
\[
\Lambda_X:\int_{\mathbb A}X\to\int_{\mathbb B}Y.
\]

\item \textbf{Identify the flattened position object.}
The flattened position object is exactly
\[
A_1\times_{t_A,A_0,p_X}X,
\]
with source \(z\) and target \(\rho_c(a,z)\).

\item \textbf{Identify the action-category span.}
These are the source and target of the action category
\[
\int_{\mathbb A}X.
\]

\item \textbf{Read the evaluation maps.}
The evaluation maps give the chosen input arrow, output arrow, downstream
transition, and continuation state, hence recover the input projection and the
arrow map of \(\Lambda_X\).
\end{enumerate}
\end{proof}

\subsection{Flattening the Mealy coalgebra}
\label{subsec:poly-relative-flattening-mealy-coalgebra}

\begin{lemma}[Flattening pullback for the Mealy coalgebra]
\label{lem:poly-relative-mealy-position-pullback}
For the relative Mealy response coalgebra
\[
c_{\Phi,Y}:S_{\Phi,Y}\to
\mathcal M_{\mathbb A,\mathbb B,Y}(S_{\Phi,Y}),
\]
the pullback of the position projection
\[
\partial_S:\mathrm{Pos}_S\to
\mathcal M_{\mathbb A,\mathbb B,Y}(S_{\Phi,Y})
\]
along \(c_{\Phi,Y}\) is canonically isomorphic to
\[
M_{\Phi,Y}
=
A_1\times_{t_A,A_0,p_P\pi_P}S_{\Phi,Y}.
\]
Under this isomorphism,
\[
\operatorname{in}=\iota_A,
\qquad
\operatorname{out}=\iota_B,
\qquad
\operatorname{down}=\lambda_{\Phi,Y},
\qquad
\operatorname{next}=t_{\Phi,Y}.
\]
\end{lemma}

\begin{proof}
\leavevmode
\begin{enumerate}
\item \textbf{Pull back positions along the coalgebra.}
Pulling back \(\partial_S\) along \(c_{\Phi,Y}\) chooses, from the total
response profile determined by a state \((x,y)\), a single input position
\[
a:u\to p_P(x).
\]

\item \textbf{Identify the resulting object.}
Thus the resulting object is precisely the object of primitive executions
\[
A_1\times_{t_A,A_0,p_P\pi_P}S_{\Phi,Y}.
\]

\item \textbf{Identify the evaluation maps.}
The four evaluation maps read off, respectively, the input arrow, the raw
emitted output, the downstream transition, and the next coupled state.  These
are exactly
\[
\iota_A,\qquad
\iota_B,\qquad
\lambda_{\Phi,Y},\qquad
t_{\Phi,Y}.
\]
\end{enumerate}
\end{proof}

\begin{theorem}[Program category recovered by flattening]
\label{thm:poly-relative-program-category-recovered}
Flattening the relative Mealy response coalgebra
\[
c_{\Phi,Y}
\]
recovers the primitive execution span
\[
\mathbf m_{\Phi,Y}
=
\left(
S_{\Phi,Y}
\xleftarrow{s_{\Phi,Y}}
M_{\Phi,Y}
\xrightarrow{t_{\Phi,Y}}
S_{\Phi,Y}
\right)
\]
of the relative program category
\[
\mathsf{Prog}_{\Phi,Y}.
\]
Explicitly,
\[
M_{\Phi,Y}
=
A_1\times_{t_A,A_0,p_P\pi_P}S_{\Phi,Y},
\]
\[
s_{\Phi,Y}(a,x,y)=(x,y),
\]
and
\[
t_{\Phi,Y}(a,x,y)
=
\left(
\delta_P(a,x),
\sigma(\omega_P(a,x),y)
\right).
\]
Consequently, the internal category recovered from the flattened coalgebra is
\[
\mathsf{Prog}_{\Phi,Y}
=
\int_{\mathbb A}\Phi^\ast Y.
\]
\end{theorem}

\begin{proof}
\leavevmode
\begin{enumerate}
\item \textbf{Use multiplicativity.}
By Theorem~\ref{thm:poly-relative-mealy-coalgebra-transpose},
\(c_{\Phi,Y}\) is multiplicative.

\item \textbf{Apply flattening.}
Theorem~\ref{thm:poly-relative-flattening-multiplicative} identifies the
flattened span with the primitive execution span of the action category.

\item \textbf{Use the Mealy pullback calculation.}
Lemma~\ref{lem:poly-relative-mealy-position-pullback} identifies this flattened
primitive object with
\[
M_{\Phi,Y}.
\]

\item \textbf{Read the explicit formulas.}
The source and target maps are the uncurried maps defining the labelled action
associated to \(\Phi\).
\end{enumerate}
\end{proof}

\begin{remark}
\label{rem:poly-relative-polynomial-not-replacement}
The polynomial coalgebra does not replace the program category.  It is the
curried total-response presentation of the same labelled one-step data whose
flattened form is the primitive execution span.
\end{remark}

\section{Output-lax profile comparisons}
\label{sec:poly-relative-output-lax-profile-comparisons}

\subsection{Why ordinary coalgebra morphisms are too strict}
\label{subsec:poly-relative-ordinary-morphisms-too-strict}

The relative response polynomial lives over
\[
J_Y=A_0\times Y.
\]
Ordinary coalgebra morphisms over \(J_Y\) are too strict for Mealy simulations,
because simulations execute output-comparison arrows in \(Y\) and therefore
change the \(Y\)-component of the interface.

Thus we isolate the correct comparison structure.

\subsection{Definition of output-lax comparison}
\label{subsec:poly-relative-output-lax-comparison-definition}

Let
\[
c:X\to\mathcal M_{\mathbb A,\mathbb B,Y}(X),
\qquad
d:Z\to\mathcal M_{\mathbb A,\mathbb B,Y}(Z)
\]
be multiplicative labelled response coalgebras.  Write their uncurried data as
\[
(\rho_X,\lambda_X)
\qquad\text{and}\qquad
(\rho_Z,\lambda_Z).
\]

\begin{definition}[Output-lax profile comparison]
\label{def:poly-relative-output-lax-profile-comparison}
An \textbf{output-lax profile comparison}
\[
(c:X\to\mathcal M(X))
\longrightarrow
(d:Z\to\mathcal M(Z))
\]
consists of maps
\[
h:X\to Z,
\qquad
\chi:X\to B_1,
\]
such that
\[
p_Zh=p_X,
\]
\[
t_B\chi=r y_X,
\qquad
s_B\chi=r y_Zh,
\]
and
\[
y_Zh=\sigma(\chi,y_X).
\]
Thus \(\chi_x\) is an arrow
\[
r(y_Zh(x))\longrightarrow r(y_Xx)
\]
in \(\mathbb B\), and it evaluates in the right action \(Y\) as an arrow
\[
y_Xx\longrightarrow y_Zh(x)
\]
of \(\int_{\mathbb B}Y\).

Writing
\[
c(x)(a)=(\beta_X(a,x),x_a)
\]
and
\[
d(hx)(a)=(\beta_Z(a,hx),z_a),
\]
the comparison also satisfies the strict input-target equation
\[
z_a=h(x_a),
\]
or equivalently
\[
\rho_Z(a,hx)=h(\rho_X(a,x)).
\]
Finally, the output-lax naturality equation is
\[
\chi_x\circ \beta_Z(a,hx)
=
\beta_X(a,x)\circ \chi_{x_a}.
\]
Equivalently,
\[
m_B(\beta_Z(a,hx),\chi_x)
=
m_B(\chi_{x_a},\beta_X(a,x)).
\]
\end{definition}

The structure equations can be displayed as
\[
\begin{tikzcd}[column sep=large,row sep=large]
&
X
  \arrow[dl,"y_X"']
  \arrow[dr,"y_Zh"]
  \arrow[d,"\chi" description]
&
\\
Y
  \arrow[d,"r"']
&
B_1
  \arrow[dl,"t_B" description]
  \arrow[dr,"s_B" description]
&
Y
  \arrow[d,"r"]
\\
B_0
&&
B_0 .
\end{tikzcd}
\]
The action condition is the additional equation
\[
y_Zh=\sigma(\chi,y_X).
\]

\subsection{Flattening output-lax comparisons}
\label{subsec:poly-relative-flatten-output-lax-comparisons}

\begin{proposition}[Flattening an output-lax profile comparison]
\label{prop:poly-relative-flatten-output-lax-profile-comparison}
An output-lax profile comparison \((h,\chi):c\to d\) induces:
\begin{enumerate}
\item an internal functor
\[
H:\int_{\mathbb A}X\to\int_{\mathbb A}Z
\]
strictly over \(\mathbb A^{\mathrm{op}}\);
\item an internal natural transformation
\[
\lambda^\chi:
\Lambda_X\Rightarrow\Lambda_ZH
\]
with codomain \(\int_{\mathbb B}Y\).
\end{enumerate}
The functor \(H\) is given by
\[
H_0=h,
\qquad
H_1(a,x)=(a,hx),
\]
and the natural transformation has component
\[
(\chi_x,y_Xx):y_Xx\to y_Zh(x)
\]
at \(x\).
\end{proposition}

\begin{proof}
\leavevmode
\begin{enumerate}
\item \textbf{Define the object map.}
Set
\[
H_0=h:X\to Z.
\]
The equation
\[
p_Zh=p_X
\]
says that this map is over \(A_0\).

\item \textbf{Define the arrow map.}
The arrow map is
\[
H_1(a,x)=(a,hx).
\]
It is well-defined because
\[
t_A(a)=p_Xx=p_Zh(x).
\]

\item \textbf{Check source preservation.}
Source preservation is immediate from
\[
s(a,x)=x
\]
and
\[
s(a,hx)=hx.
\]

\item \textbf{Check target preservation.}
The strict input-target equation
\[
\rho_Z(a,hx)=h(\rho_X(a,x))
\]
says exactly that \(H\) preserves targets of primitive arrows.

\item \textbf{Check identities and composition.}
Since the input arrow \(a\) is unchanged and \(h\) preserves the action targets,
identity and composition preservation follow from multiplicativity of the two
coalgebras.

\item \textbf{Construct the downstream natural transformation.}
The component
\[
(\chi_x,y_Xx)
\]
is a well-defined arrow of \(\int_{\mathbb B}Y\) because
\[
t_B\chi_x=r(y_Xx)
\]
and
\[
y_Zh(x)=\sigma(\chi_x,y_Xx).
\]

\item \textbf{Check naturality.}
The output-lax naturality equation
\[
\chi_x\circ \beta_Z(a,hx)
=
\beta_X(a,x)\circ \chi_{x_a}
\]
is exactly the naturality square in the downstream action category.
\end{enumerate}
\end{proof}

The naturality square is
\[
\begin{tikzcd}[column sep=large,row sep=large]
y_Xx
  \arrow[r,"{(\beta_X(a,x),y_Xx)}"]
  \arrow[d,"{(\chi_x,y_Xx)}"']
&
y_Xx_a
  \arrow[d,"{(\chi_{x_a},y_Xx_a)}"]
\\
y_Zh(x)
  \arrow[r,"{(\beta_Z(a,hx),y_Zh(x))}"']
&
y_Zh(x_a).
\end{tikzcd}
\]

\section{Mealy simulations as output-lax profile comparisons}
\label{sec:poly-relative-mealy-simulations-profile-comparisons}

\subsection{Evaluated comparison data}
\label{subsec:poly-relative-simulation-evaluated-data}

Let
\[
\Phi=(P,\delta_P,\omega_P),
\qquad
\Psi=(Q,\delta_Q,\omega_Q)
:
\mathbb A\rightsquigarrow\mathbb B
\]
be Mealy morphisms, and let
\[
\Theta:\Phi\Rightarrow\Psi
\]
be a Mealy simulation with components
\[
\theta:Q\to P,
\qquad
\alpha:Q\to B_1.
\]
Thus
\[
p_P\theta=p_Q,
\qquad
s_B\alpha=q_P\theta,
\qquad
t_B\alpha=q_Q,
\]
and
\[
\theta(\delta_Q(a,y))
=
\delta_P(a,\theta y),
\]
\[
\alpha_y\circ\omega_P(a,\theta y)
=
\omega_Q(a,y)\circ\alpha_{\delta_Q(a,y)}.
\]

Let \(Y\) be a right \(\mathbb B\)-action.  The evaluated state comparison is
\[
h_{\Theta,Y}:S_{\Psi,Y}\to S_{\Phi,Y},
\qquad
h_{\Theta,Y}(y,z)
=
\bigl(\theta y,\sigma(\alpha_y,z)\bigr).
\]
Define
\[
\chi_{\Theta,Y}:S_{\Psi,Y}\to B_1
\]
by
\[
\chi_{\Theta,Y}(y,z)=\alpha_y.
\]

\subsection{Simulations induce output-lax comparisons}
\label{subsec:poly-relative-simulation-output-lax-theorem}

\begin{theorem}[Simulations induce output-lax profile comparisons]
\label{thm:poly-relative-simulations-profile-comparisons}
The pair
\[
(h_{\Theta,Y},\chi_{\Theta,Y})
\]
is an output-lax profile comparison
\[
c_{\Psi,Y}\longrightarrow c_{\Phi,Y}.
\]
After flattening, it recovers the input-strict functor
\[
H_{\Theta,Y}:
\mathsf{Prog}_{\Psi,Y}\to\mathsf{Prog}_{\Phi,Y}
\]
and the downstream comparison natural transformation
\[
\lambda_{\Theta,Y}:
\Lambda_{\Psi,Y}
\Rightarrow
\Lambda_{\Phi,Y}H_{\Theta,Y}.
\]
\end{theorem}

\begin{proof}
\leavevmode
\begin{enumerate}
\item \textbf{Check input strictness.}
The map \(h_{\Theta,Y}\) is strict over \(A_0\) because
\[
p_P\theta=p_Q.
\]

\item \textbf{Check the output-comparison typing.}
The map \(\chi_{\Theta,Y}\) satisfies
\[
t_B\chi_{\Theta,Y}(y,z)=t_B\alpha_y=q_Qy=r(z),
\]
and
\[
s_B\chi_{\Theta,Y}(y,z)=s_B\alpha_y=q_P\theta y
=
r(\sigma(\alpha_y,z)).
\]

\item \textbf{Check evaluation in the downstream action.}
Moreover,
\[
\pi_Yh_{\Theta,Y}(y,z)=\sigma(\alpha_y,z)
=
\sigma(\chi_{\Theta,Y}(y,z),z).
\]

\item \textbf{Compute the source response.}
The source coalgebra \(c_{\Psi,Y}\) has
\[
c_{\Psi,Y}(y,z)(a)
=
\left(
\omega_Q(a,y),
\left(
\delta_Q(a,y),
\sigma(\omega_Q(a,y),z)
\right)
\right).
\]

\item \textbf{Compute the target response at the compared state.}
The target coalgebra at \(h_{\Theta,Y}(y,z)\) has
\[
c_{\Phi,Y}(\theta y,\sigma(\alpha_y,z))(a)
=
\left(
\omega_P(a,\theta y),
\left(
\delta_P(a,\theta y),
\sigma(\omega_P(a,\theta y),\sigma(\alpha_y,z))
\right)
\right).
\]

\item \textbf{Check strict target preservation.}
The strict input-target equation is
\[
\rho_{\Phi}(a,h_{\Theta,Y}(y,z))
=
h_{\Theta,Y}(\rho_{\Psi}(a,y,z)).
\]
Indeed, the left-hand side is
\[
\left(
\delta_P(a,\theta y),
\sigma(\omega_P(a,\theta y),\sigma(\alpha_y,z))
\right),
\]
while the right-hand side is
\[
\left(
\theta\delta_Q(a,y),
\sigma(\alpha_{\delta_Q(a,y)},\sigma(\omega_Q(a,y),z))
\right).
\]
The first components agree by the state part of the simulation law.  The second
components agree by the output part of the simulation law and the right-action
multiplication law for \(Y\).

\item \textbf{Check output-lax naturality.}
The output-lax square is precisely
\[
\alpha_y\circ\omega_P(a,\theta y)
=
\omega_Q(a,y)\circ\alpha_{\delta_Q(a,y)}.
\]
Therefore \((h_{\Theta,Y},\chi_{\Theta,Y})\) is an output-lax profile
comparison.

\item \textbf{Flatten.}
The final statement is
Proposition~\ref{prop:poly-relative-flatten-output-lax-profile-comparison}
specialized to the Mealy-induced coalgebras.
\end{enumerate}
\end{proof}

\begin{remark}
\label{rem:poly-relative-not-ordinary-coalgebra-map}
The map
\[
h_{\Theta,Y}:S_{\Psi,Y}\to S_{\Phi,Y}
\]
is generally not a map over \(J_Y=A_0\times Y\), since it changes the
downstream component from \(z\) to \(\sigma(\alpha_y,z)\).  The correct
polynomial comparison is therefore output-lax, not an ordinary coalgebra
morphism in \(\mathcal E/J_Y\).
\end{remark}

\section{Guard loci and guarded primitives}
\label{sec:poly-relative-guard-loci-primitives}

\subsection{Guard loci}
\label{subsec:poly-relative-guard-loci}

Fix \(\Phi\) and \(Y\), and write
\[
S:=S_{\Phi,Y},
\qquad
M:=M_{\Phi,Y}.
\]

A guarded primitive specification may involve the following predicates:
\[
\theta\hookrightarrow S
\qquad
\text{current coupled-state guard},
\]
\[
R\hookrightarrow A_1
\qquad
\text{input-arrow guard},
\]
\[
O\hookrightarrow B_1
\qquad
\text{raw output-arrow guard},
\]
\[
T\hookrightarrow\left(\int_{\mathbb B}Y\right)_1
\qquad
\text{downstream transition guard},
\]
and
\[
\varphi\hookrightarrow S
\qquad
\text{continuation-state predicate}.
\]

The primitive execution object has canonical maps
\[
\iota_A:M\to A_1,
\qquad
\iota_A(a,x,y)=a,
\]
\[
\iota_B:M\to B_1,
\qquad
\iota_B(a,x,y)=\omega_P(a,x),
\]
and
\[
\lambda_{\Phi,Y}:M\to\left(\int_{\mathbb B}Y\right)_1,
\qquad
\lambda_{\Phi,Y}(a,x,y)=(\omega_P(a,x),y).
\]
They are displayed by
\[
\begin{tikzcd}[column sep=large,row sep=large]
&
M
  \arrow[dl,"\iota_A"']
  \arrow[d,"\lambda_{\Phi,Y}" description]
  \arrow[dr,"\iota_B"]
&
\\
A_1
&
\left(\int_{\mathbb B}Y\right)_1
  \arrow[l,phantom,"\mapsto"]
  \arrow[r,"\pi_{B_1}"']
&
B_1 .
\end{tikzcd}
\]

In the terminal downstream case,
\[
\int_{\mathbb B}\mathbf 1_{\mathbb B}\cong\mathbb B^{\mathrm{op}},
\]
so downstream transition guards reduce to predicates on output arrows of
\(\mathbb B\).

\subsection{Guarded primitive programs}
\label{subsec:poly-relative-guarded-primitives}

Let
\[
\theta\hookrightarrow S,
\qquad
R\hookrightarrow A_1,
\qquad
O\hookrightarrow B_1,
\qquad
T\hookrightarrow\left(\int_{\mathbb B}Y\right)_1
\]
be guards.

Pull them back to \(M\):
\[
\theta^\sharp:=s_{\Phi,Y}^\ast\theta,
\]
\[
R^\sharp:=\iota_A^\ast R,
\qquad
O^\sharp:=\iota_B^\ast O,
\qquad
T^\sharp:=\lambda_{\Phi,Y}^\ast T.
\]
Their simultaneous restriction is
\[
(\theta;R/O/T)^\sharp
:=
\theta^\sharp\wedge R^\sharp\wedge O^\sharp\wedge T^\sharp
\hookrightarrow M.
\]

\begin{definition}[Guarded primitive relative program]
\label{def:poly-relative-guarded-primitive}
The \textbf{guarded primitive relative program}
\[
\theta;R/O/T
\]
is the restricted span
\[
S
\xleftarrow{s_{\Phi,Y}i}
(\theta;R/O/T)^\sharp
\xrightarrow{t_{\Phi,Y}i}
S,
\]
where
\[
i:(\theta;R/O/T)^\sharp\hookrightarrow M
\]
is the inclusion.  Its diamond is denoted
\[
\langle\theta;R/O/T\rangle\varphi.
\]
\end{definition}

The restricted span is displayed by
\[
\begin{tikzcd}[column sep=large,row sep=large]
&
(\theta;R/O/T)^\sharp
  \arrow[dl,"s_{\Phi,Y}i"']
  \arrow[dr,"t_{\Phi,Y}i"]
  \arrow[d,hook,"i" description]
&
\\
S
&
M
  \arrow[l,"s_{\Phi,Y}"]
  \arrow[r,"t_{\Phi,Y}"']
&
S .
\end{tikzcd}
\]

\begin{proposition}[External reading of guarded primitives]
\label{prop:poly-relative-guarded-primitive-external}
In \(\mathbf{Set}\),
\[
(x,y)\models\langle\theta;R/O/T\rangle\varphi
\]
holds precisely when
\[
\theta(x,y)
\]
and there exists an admissible input arrow
\[
a:u\to p_P(x)
\]
such that
\[
R(a),
\qquad
O(\omega_P(a,x)),
\qquad
T(\omega_P(a,x),y),
\]
and
\[
\left(
\delta_P(a,x),
\sigma(\omega_P(a,x),y)
\right)
\models\varphi.
\]
\end{proposition}

\begin{proof}
\leavevmode
\begin{enumerate}
\item \textbf{Expand the restricted apex.}
The restricted primitive span has apex the subobject of \(M\) consisting of
triples \((a,x,y)\) satisfying the four pulled-back guards.

\item \textbf{Expand the diamond.}
The diamond says that there exists such a triple with source \((x,y)\) and
target satisfying \(\varphi\).

\item \textbf{Read the target.}
The target is
\[
\bigl(
\delta_P(a,x),
\sigma(\omega_P(a,x),y)
\bigr).
\]
\end{enumerate}
\end{proof}

\section{Guarded polynomial response liftings}
\label{sec:poly-relative-guarded-response-liftings}

\subsection{Existential profile liftings}
\label{subsec:poly-relative-existential-profile-liftings}

Assume \(\mathcal E\) is regular.  Let
\[
c_{\Phi,Y}:
S\to\mathcal M_{\mathbb A,\mathbb B,Y}(S)
\]
be the relative response coalgebra, and let
\[
\partial_S:\mathrm{Pos}_S\to\mathcal M_{\mathbb A,\mathbb B,Y}(S)
\]
be the position projection.

The position object carries evaluation maps
\[
\operatorname{in}:\mathrm{Pos}_S\to A_1,
\]
\[
\operatorname{out}:\mathrm{Pos}_S\to B_1,
\]
\[
\operatorname{down}:\mathrm{Pos}_S\to
\left(\int_{\mathbb B}Y\right)_1,
\]
and
\[
\operatorname{next}:\mathrm{Pos}_S\to S.
\]

The following liftings are coalgebraic predicate liftings for the response
polynomial, later pulled back along the coalgebra to the span modalities of
the flattened program category.  This is the same general pattern by which
modal operators for coalgebras are described using predicate liftings
\cite{Pattinson2003CoalgebraicModalLogic}.

\begin{definition}[Guarded existential profile lifting]
\label{def:poly-relative-guarded-existential-lifting}
Let
\[
R\hookrightarrow A_1,
\qquad
O\hookrightarrow B_1,
\qquad
T\hookrightarrow\left(\int_{\mathbb B}Y\right)_1,
\qquad
\varphi\hookrightarrow S
\]
be predicates.  The \textbf{guarded existential profile lifting} is the
predicate
\[
\lambda_{R/O/T}(\varphi)
\hookrightarrow
\mathcal M_{\mathbb A,\mathbb B,Y}(S)
\]
defined by
\[
\lambda_{R/O/T}(\varphi)
:=
\exists_{\partial_S}
\left(
\operatorname{in}^{\ast}R
\wedge
\operatorname{out}^{\ast}O
\wedge
\operatorname{down}^{\ast}T
\wedge
\operatorname{next}^{\ast}\varphi
\right).
\]
\end{definition}

Externally, a response profile satisfies
\[
\lambda_{R/O/T}(\varphi)
\]
when some input position satisfies \(R\), the selected raw output satisfies
\(O\), the downstream transition satisfies \(T\), and the continuation state
satisfies \(\varphi\).

\subsection{Agreement with span diamonds}
\label{subsec:poly-relative-lifting-span-agreement}

\begin{theorem}[Agreement of profile liftings and span diamonds]
\label{thm:poly-relative-lifting-span-agreement}
For every state guard
\[
\theta\hookrightarrow S,
\]
input guard
\[
R\hookrightarrow A_1,
\]
output guard
\[
O\hookrightarrow B_1,
\]
downstream transition guard
\[
T\hookrightarrow\left(\int_{\mathbb B}Y\right)_1,
\]
and postcondition
\[
\varphi\hookrightarrow S,
\]
one has
\[
\theta\wedge c_{\Phi,Y}^{\ast}\lambda_{R/O/T}(\varphi)
=
\langle\theta;R/O/T\rangle\varphi.
\]
\end{theorem}

\begin{proof}
\leavevmode
\begin{enumerate}
\item \textbf{Pull back the position projection.}
By Lemma~\ref{lem:poly-relative-mealy-position-pullback}, pulling back the
position projection
\[
\partial_S:\mathrm{Pos}_S\to\mathcal M_{\mathbb A,\mathbb B,Y}(S)
\]
along \(c_{\Phi,Y}\) gives the primitive execution object \(M\).

\item \textbf{Identify the evaluation maps.}
Under this identification,
\[
\operatorname{in}=\iota_A,
\qquad
\operatorname{out}=\iota_B,
\qquad
\operatorname{down}=\lambda_{\Phi,Y},
\qquad
\operatorname{next}=t_{\Phi,Y}.
\]

\item \textbf{Apply Beck--Chevalley.}
By Beck--Chevalley for regular images,
\[
c_{\Phi,Y}^{\ast}\lambda_{R/O/T}(\varphi)
=
\exists_{s_{\Phi,Y}}
\left(
\iota_A^\ast R
\wedge
\iota_B^\ast O
\wedge
\lambda_{\Phi,Y}^\ast T
\wedge
t_{\Phi,Y}^\ast\varphi
\right).
\]

\item \textbf{Pull in the state guard.}
Using Frobenius reciprocity,
\[
\theta\wedge c_{\Phi,Y}^{\ast}\lambda_{R/O/T}(\varphi)
=
\exists_{s_{\Phi,Y}}
\left(
s_{\Phi,Y}^\ast\theta
\wedge
\iota_A^\ast R
\wedge
\iota_B^\ast O
\wedge
\lambda_{\Phi,Y}^\ast T
\wedge
t_{\Phi,Y}^\ast\varphi
\right).
\]

\item \textbf{Identify the restricted diamond.}
The right-hand side is exactly the diamond of the restricted primitive span
\[
\theta;R/O/T.
\]
\end{enumerate}
\end{proof}

\begin{remark}
\label{rem:poly-relative-central-comparison}
The polynomial lifting is a predicate on total response profiles.  Pulling it
back along the relative Mealy coalgebra recovers the guarded span diamond on
the flattened program category.
\end{remark}

\section{Profile boxes}
\label{sec:poly-relative-profile-boxes}

\subsection{Weak and strong boxes}
\label{subsec:poly-relative-weak-strong-boxes}

Assume \(\mathcal E\) is a first-order Heyting category.

Let
\[
R\hookrightarrow A_1,
\qquad
O\hookrightarrow B_1,
\qquad
T\hookrightarrow\left(\int_{\mathbb B}Y\right)_1,
\qquad
\varphi\hookrightarrow S
\]
be predicates.

\begin{definition}[Weak universal profile box]
\label{def:poly-relative-weak-profile-box}
The \textbf{weak universal profile box} is
\[
\Box^{\mathrm{prof}}_{R/O/T}(\varphi)
:=
\forall_{\partial_S}
\left(
(\operatorname{in}^{\ast}R
\wedge
\operatorname{out}^{\ast}O
\wedge
\operatorname{down}^{\ast}T)
\Rightarrow
\operatorname{next}^{\ast}\varphi
\right).
\]
\end{definition}

Externally, this says that every position whose input, output, and downstream
transition satisfy the guards has a continuation satisfying \(\varphi\).

\begin{definition}[Strong universal profile box]
\label{def:poly-relative-strong-profile-box}
The \textbf{strong universal profile box} is
\[
\Box^{\mathrm{prof,strong}}_{R/O/T}(\varphi)
:=
\forall_{\partial_S}
\left(
\operatorname{in}^{\ast}R
\Rightarrow
\left(
\operatorname{out}^{\ast}O
\wedge
\operatorname{down}^{\ast}T
\wedge
\operatorname{next}^{\ast}\varphi
\right)
\right).
\]
\end{definition}

Externally, this says that every \(R\)-input position produces an output
satisfying \(O\), a downstream transition satisfying \(T\), and a continuation
satisfying \(\varphi\).

\subsection{Agreement with span boxes}
\label{subsec:poly-relative-profile-box-span-agreement}

\begin{theorem}[Agreement of weak profile boxes and span boxes]
\label{thm:poly-relative-profile-box-span-agreement}
One has
\[
c_{\Phi,Y}^{\ast}
\Box^{\mathrm{prof}}_{R/O/T}(\varphi)
=
[R/O/T]\varphi,
\]
where \([R/O/T]\varphi\) denotes the box modality of the restricted primitive
span determined by the guards \(R,O,T\).

More generally, for a state guard
\[
\theta\hookrightarrow S,
\]
one has
\[
[\theta;R/O/T]\varphi
=
\theta\Rightarrow
\bigl(
c_{\Phi,Y}^{\ast}
\Box^{\mathrm{prof}}_{R/O/T}(\varphi)
\bigr).
\]
\end{theorem}

\begin{proof}
\leavevmode
\begin{enumerate}
\item \textbf{Apply Beck--Chevalley for \(\forall\).}
By Beck--Chevalley for \(\forall\), pulling back
\(\Box^{\mathrm{prof}}_{R/O/T}(\varphi)\) along \(c_{\Phi,Y}\) gives
\[
\forall_{s_{\Phi,Y}}
\left(
(\iota_A^\ast R
\wedge
\iota_B^\ast O
\wedge
\lambda_{\Phi,Y}^\ast T)
\Rightarrow
t_{\Phi,Y}^\ast\varphi
\right).
\]

\item \textbf{Identify the restricted box.}
This is exactly the box of the restricted primitive span determined by
\(R,O,T\).

\item \textbf{Add the state guard.}
For the guarded version, use the test-prefix box law:
\[
[\theta?;U]\varphi
=
\theta\Rightarrow[U]\varphi.
\]
\end{enumerate}
\end{proof}

\begin{remark}[Weak versus strong profile boxes]
\label{rem:poly-relative-weak-strong-profile-boxes}
The weak profile box pulls back to the usual span box for the restricted
primitive span.  The strong profile box is generally stricter: it requires
every \(R\)-input position to produce an output satisfying \(O\), a downstream
transition satisfying \(T\), and a continuation satisfying \(\varphi\).
\end{remark}

\section{Simulation stability after flattening}
\label{sec:poly-relative-simulation-stability}

\subsection{Guard transport along simulations}
\label{subsec:poly-relative-guard-transport-simulation}

Let
\[
\Theta:\Phi\Rightarrow\Psi
\]
be a Mealy simulation, and let \(Y\) be a right \(\mathbb B\)-action.  The
induced functor
\[
H_{\Theta,Y}:
\mathsf{Prog}_{\Psi,Y}\to\mathsf{Prog}_{\Phi,Y}
\]
has object map
\[
h_{\Theta,Y}:S_{\Psi,Y}\to S_{\Phi,Y}
\]
and arrow map
\[
k_{\Theta,Y}:M_{\Psi,Y}\to M_{\Phi,Y},
\]
given by
\[
k_{\Theta,Y}(a,y,z)
=
\bigl(a,\theta y,\sigma(\alpha_y,z)\bigr).
\]

The source square is a pullback:
\[
\begin{tikzcd}[column sep=large,row sep=large]
M_{\Psi,Y}
  \arrow[r,"k_{\Theta,Y}"]
  \arrow[d,"s_\Psi"']
&
M_{\Phi,Y}
  \arrow[d,"s_\Phi"]
\\
S_{\Psi,Y}
  \arrow[r,"h_{\Theta,Y}"']
&
S_{\Phi,Y}
\arrow[ul,phantom,very near start,"\lrcorner"] .
\end{tikzcd}
\]

\begin{proposition}[Guard transport along simulations]
\label{prop:poly-relative-guard-transport-simulation}
Let
\[
U\hookrightarrow M_{\Phi,Y}
\]
be a primitive subprogram for \(\Phi\).  Its pullback along the simulation is
\[
k_{\Theta,Y}^\ast U\hookrightarrow M_{\Psi,Y}.
\]
If \(U\) is presented by guards \(R,O,T\) on the target side, then
\(k_{\Theta,Y}^\ast U\) is the source-side predicate requiring
\[
R(a),
\]
\[
O(\omega_P(a,\theta y)),
\]
and
\[
T(\omega_P(a,\theta y),\sigma(\alpha_y,z)).
\]
If the target-side subprogram also has a state guard
\[
\theta_0\hookrightarrow S_{\Phi,Y},
\]
then its pullback additionally requires
\[
\theta_0(\theta y,\sigma(\alpha_y,z)).
\]
Equivalently, the state guard is pulled back as
\[
h_{\Theta,Y}^{\ast}\theta_0.
\]
\end{proposition}

\begin{proof}
\leavevmode
\begin{enumerate}
\item \textbf{Track the input component.}
The map \(k_{\Theta,Y}\) preserves the input component \(a\).  Hence input
guards pull back by direct substitution:
\[
R(a).
\]

\item \textbf{Track the target output component.}
The output and downstream transition components are those of the
\(\Phi\)-profile at the evaluated state
\[
(\theta y,\sigma(\alpha_y,z)).
\]
Thus the output guard becomes
\[
O(\omega_P(a,\theta y)).
\]

\item \textbf{Track the downstream transition.}
The downstream transition guard becomes
\[
T(\omega_P(a,\theta y),\sigma(\alpha_y,z)).
\]

\item \textbf{Track state guards.}
A state guard is pulled back along the object map \(h_{\Theta,Y}\), hence
becomes
\[
h_{\Theta,Y}^{\ast}\theta_0.
\]
\end{enumerate}
\end{proof}

\subsection{Modal stability under simulations}
\label{subsec:poly-relative-modal-stability-simulation}

\begin{proposition}[Diamond stability under simulations]
\label{prop:poly-relative-diamond-stability-simulation}
Assume \(\mathcal E\) is regular.  For every primitive subprogram
\[
U\hookrightarrow M_{\Phi,Y}
\]
and every predicate
\[
\varphi\hookrightarrow S_{\Phi,Y},
\]
one has
\[
h_{\Theta,Y}^\ast\langle U\rangle\varphi
=
\langle k_{\Theta,Y}^\ast U\rangle h_{\Theta,Y}^\ast\varphi.
\]
\end{proposition}

\begin{proof}
\leavevmode
\begin{enumerate}
\item \textbf{Pull back the subprogram.}
Let
\[
i:U\hookrightarrow M_{\Phi,Y}
\]
be the subprogram inclusion, and form the pullback
\[
\begin{tikzcd}[column sep=large,row sep=large]
k_{\Theta,Y}^\ast U
  \arrow[r,"\bar k_{\Theta,Y}"]
  \arrow[d,"\bar i"']
&
U
  \arrow[d,"i"]
\\
M_{\Psi,Y}
  \arrow[r,"k_{\Theta,Y}"']
&
M_{\Phi,Y}
\arrow[ul,phantom,very near start,"\lrcorner"] .
\end{tikzcd}
\]

\item \textbf{Use the source pullback square.}
Since the square
\[
\begin{tikzcd}[column sep=large,row sep=large]
M_{\Psi,Y} \arrow[r,"k_{\Theta,Y}"] \arrow[d,"s_\Psi"']
&
M_{\Phi,Y} \arrow[d,"s_\Phi"]
\\
S_{\Psi,Y} \arrow[r,"h_{\Theta,Y}"']
&
S_{\Phi,Y}
\end{tikzcd}
\]
is a pullback, the induced square
\[
\begin{tikzcd}[column sep=large,row sep=large]
k_{\Theta,Y}^\ast U
  \arrow[r,"\bar k_{\Theta,Y}"]
  \arrow[d,"s_\Psi\bar i"']
&
U
  \arrow[d,"s_\Phi i"]
\\
S_{\Psi,Y}
  \arrow[r,"h_{\Theta,Y}"']
&
S_{\Phi,Y}
\arrow[ul,phantom,very near start,"\lrcorner"]
\end{tikzcd}
\]
is also a pullback.

\item \textbf{Use target compatibility.}
Since \(H_{\Theta,Y}\) is an internal functor,
\[
t_\Phi i\bar k_{\Theta,Y}=h_{\Theta,Y}t_\Psi\bar i.
\]

\item \textbf{Apply Beck--Chevalley.}
Beck--Chevalley for regular images gives
\[
h_{\Theta,Y}^\ast
\exists_{s_\Phi i}\bigl((t_\Phi i)^\ast\varphi\bigr)
=
\exists_{s_\Psi\bar i}
\bigl((t_\Psi\bar i)^\ast h_{\Theta,Y}^\ast\varphi\bigr).
\]

\item \textbf{Identify the two diamonds.}
The left-hand side is
\[
h_{\Theta,Y}^\ast\langle U\rangle\varphi,
\]
and the right-hand side is
\[
\langle k_{\Theta,Y}^\ast U\rangle h_{\Theta,Y}^\ast\varphi.
\]
\end{enumerate}
\end{proof}

\begin{proposition}[Box stability under simulations]
\label{prop:poly-relative-box-stability-simulation}
Assume \(\mathcal E\) is first-order Heyting.  For every primitive subprogram
\[
U\hookrightarrow M_{\Phi,Y}
\]
and every predicate
\[
\varphi\hookrightarrow S_{\Phi,Y},
\]
one has
\[
h_{\Theta,Y}^\ast[U]\varphi
=
[k_{\Theta,Y}^\ast U]h_{\Theta,Y}^\ast\varphi.
\]
\end{proposition}

\begin{proof}
\leavevmode
\begin{enumerate}
\item \textbf{Use the same pullback square.}
Use the pullback square for
\[
k_{\Theta,Y}^\ast U
\]
constructed in the proof of
Proposition~\ref{prop:poly-relative-diamond-stability-simulation}.

\item \textbf{Use target compatibility.}
Again,
\[
t_\Phi i\bar k_{\Theta,Y}=h_{\Theta,Y}t_\Psi\bar i.
\]

\item \textbf{Apply Beck--Chevalley for \(\forall\).}
Beck--Chevalley for the right adjoints \(\forall\) gives the displayed
equality.
\end{enumerate}
\end{proof}

\begin{remark}
\label{rem:poly-relative-simulation-stability-meaning}
Because simulations are output-lax, guard transport is not obtained by
substituting the source output \(\omega_Q\) for the target output
\(\omega_P\).  It is obtained by pulling the target subprogram back along the
actual flattened functor \(H_{\Theta,Y}\).
\end{remark}

\section{Tests, guarded commands, and finite choice}
\label{sec:poly-relative-tests-choice}

\subsection{Tests and guarded commands}
\label{subsec:poly-relative-tests-guarded-commands}

The remaining guarded program operations are inherited from the span dynamic
logic on the relative state object
\[
S=S_{\Phi,Y}.
\]

Assume \(\mathcal E\) is regular.  For a predicate
\[
\theta:\Theta\hookrightarrow S,
\]
the corresponding test span is
\[
\theta?
=
(S\xleftarrow{\theta}\Theta\xrightarrow{\theta}S).
\]
It is displayed by
\[
\begin{tikzcd}
&
\Theta
  \arrow[dl,"\theta"']
  \arrow[dr,"\theta"]
&
\\
S && S .
\end{tikzcd}
\]

Let
\[
U:S\nrightarrow S
\]
be an endoprogram.  The guarded command
\[
\theta?;U
\]
is the sequential composite of the test and \(U\).

\begin{proposition}[Guard-prefix law]
\label{prop:poly-relative-guard-prefix-law}
For every endoprogram \(U:S\nrightarrow S\),
\[
\langle\theta?;U\rangle\varphi
=
\theta\wedge\langle U\rangle\varphi.
\]
\end{proposition}

\begin{proof}
\leavevmode
\begin{enumerate}
\item \textbf{Use diamond functoriality.}
By functoriality of diamonds,
\[
\langle\theta?;U\rangle\varphi
=
\langle\theta?\rangle(\langle U\rangle\varphi).
\]

\item \textbf{Use the test law.}
The test law gives
\[
\langle\theta?\rangle\psi=\theta\wedge\psi.
\]

\item \textbf{Substitute.}
Substituting
\[
\psi=\langle U\rangle\varphi
\]
gives the claim.
\end{enumerate}
\end{proof}

\subsection{Finite guarded choice}
\label{subsec:poly-relative-finite-guarded-choice}

Assume now that \(\mathcal E\) is coherent.  Let
\[
T:S\nrightarrow S
\]
be a fixed ambient endospan, and let
\[
U,V\hookrightarrow T
\]
be subprograms.  If
\[
\theta,\epsilon\hookrightarrow S
\]
are state guards, then we identify
\[
\theta?;U
\]
with the subobject
\[
s_T^\ast\theta\wedge U\hookrightarrow T,
\]
and similarly identify
\[
\epsilon?;V
\]
with
\[
s_T^\ast\epsilon\wedge V\hookrightarrow T.
\]
Their guarded choice is the join of these subobjects in \(\mathrm{Sub}(T)\):
\[
(\theta?;U)\vee(\epsilon?;V)\hookrightarrow T.
\]

\begin{proposition}[Diamond law for guarded choice]
\label{prop:poly-relative-guarded-choice-diamond}
For guarded branches inside a common ambient endospan,
\[
\left\langle
(\theta?;U)\vee(\epsilon?;V)
\right\rangle\varphi
=
(\theta\wedge\langle U\rangle\varphi)
\vee
(\epsilon\wedge\langle V\rangle\varphi).
\]
\end{proposition}

\begin{proof}
\leavevmode
\begin{enumerate}
\item \textbf{Apply finite choice.}
By the finite choice law for subprograms in a coherent category,
\[
\left\langle
(\theta?;U)\vee(\epsilon?;V)
\right\rangle\varphi
=
\langle\theta?;U\rangle\varphi
\vee
\langle\epsilon?;V\rangle\varphi.
\]

\item \textbf{Apply the guard-prefix law.}
Using Proposition~\ref{prop:poly-relative-guard-prefix-law} on both branches
gives the displayed formula.
\end{enumerate}
\end{proof}

\subsection{Internal conditionals}
\label{subsec:poly-relative-conditionals}

\begin{definition}[Internal conditional]
\label{def:poly-relative-conditional}
Let
\[
T:S\nrightarrow S
\]
be a fixed ambient endospan, and let
\[
U,V\hookrightarrow T
\]
be subprograms.  Let
\[
\theta,\epsilon\hookrightarrow S
\]
be complementary guards, meaning
\[
\theta\wedge\epsilon=\bot,
\qquad
\theta\vee\epsilon=\top.
\]
The \textbf{internal conditional} is
\[
\mathbf{if}\ \theta\ \mathbf{then}\ U\ \mathbf{else}\ V
:=
(\theta?;U)\vee(\epsilon?;V)
\hookrightarrow T.
\]
In Boolean fibres one may take
\[
\epsilon=\neg\theta.
\]
\end{definition}

\section{Boxes for guarded choice}
\label{sec:poly-relative-boxes-choice}

\subsection{Boxes of tests and guarded commands}
\label{subsec:poly-relative-boxes-tests-guarded}

Assume \(\mathcal E\) is a first-order Heyting category.

For a test,
\[
[\theta?]\varphi=\theta\Rightarrow\varphi.
\]
Therefore
\[
[\theta?;U]\varphi
=
\theta\Rightarrow[U]\varphi.
\]

\subsection{Boxes of finite choice}
\label{subsec:poly-relative-boxes-finite-choice}

\begin{proposition}[Box of finite subprogram choice]
\label{prop:poly-relative-box-finite-choice}
For subprograms
\[
W,Z\hookrightarrow T
\]
inside a fixed ambient endospan \(T:S\nrightarrow S\),
\[
[W\vee Z]\varphi
=
[W]\varphi\wedge[Z]\varphi.
\]
\end{proposition}

\begin{proof}
\leavevmode
\begin{enumerate}
\item \textbf{Write the ambient span.}
Write
\[
T=(S\xleftarrow{\ell_T}T\xrightarrow{r_T}S).
\]

\item \textbf{Rewrite restricted boxes in the ambient fibre.}
For a restricted subspan \(i_W:W\hookrightarrow T\),
\[
[W]\varphi
=
\forall_{\ell_T i_W}\bigl((r_T i_W)^\ast\varphi\bigr).
\]
Equivalently,
\[
[W]\varphi
=
\forall_{\ell_T}\bigl(W\Rightarrow r_T^\ast\varphi\bigr),
\]
using the adjunction for the monomorphism \(i_W:W\hookrightarrow T\).

\item \textbf{Apply the same formula to the join.}
Hence
\[
[W\vee Z]\varphi
=
\forall_{\ell_T}
\bigl((W\vee Z)\Rightarrow r_T^\ast\varphi\bigr).
\]

\item \textbf{Use Heyting distributivity.}
In a Heyting algebra,
\[
(W\vee Z)\Rightarrow A
=
(W\Rightarrow A)\wedge(Z\Rightarrow A).
\]

\item \textbf{Use meet preservation by \(\forall\).}
Since \(\forall_{\ell_T}\) is a right adjoint, it preserves finite meets.
Therefore
\[
[W\vee Z]\varphi
=
[W]\varphi\wedge[Z]\varphi.
\]
\end{enumerate}
\end{proof}

\subsection{Boxes for conditionals}
\label{subsec:poly-relative-boxes-conditionals}

\begin{proposition}[Box law for conditionals]
\label{prop:poly-relative-conditional-box-law}
For complementary guards \(\theta,\epsilon\) and subprograms
\[
U,V\hookrightarrow T
\]
inside a common ambient endospan,
\[
\left[
\mathbf{if}\ \theta\ \mathbf{then}\ U\ \mathbf{else}\ V
\right]\varphi
=
(\theta\Rightarrow[U]\varphi)
\wedge
(\epsilon\Rightarrow[V]\varphi).
\]
\end{proposition}

\begin{proof}
\leavevmode
\begin{enumerate}
\item \textbf{Expand the conditional.}
The conditional is
\[
(\theta?;U)\vee(\epsilon?;V).
\]

\item \textbf{Apply the box law for finite choice.}
Proposition~\ref{prop:poly-relative-box-finite-choice} gives
\[
[(\theta?;U)\vee(\epsilon?;V)]\varphi
=
[\theta?;U]\varphi\wedge[\epsilon?;V]\varphi.
\]

\item \textbf{Use the test-prefix box law.}
The test-prefix box law gives
\[
[\theta?;U]\varphi=\theta\Rightarrow[U]\varphi
\]
and
\[
[\epsilon?;V]\varphi=\epsilon\Rightarrow[V]\varphi.
\]
\end{enumerate}
\end{proof}

\section{Domain and dynamic tests}
\label{sec:poly-relative-domain-tests}

\subsection{Domain predicates}
\label{subsec:poly-relative-domain-predicate}

Let
\[
U:S\nrightarrow S
\]
be an endoprogram.

\begin{definition}[Domain predicate]
\label{def:poly-relative-domain}
The \textbf{domain predicate} of \(U\) is
\[
\operatorname{dom}(U):=\langle U\rangle\top.
\]
\end{definition}

Thus \(\operatorname{dom}(U)\) is the predicate of states from which some
\(U\)-execution is possible.

\begin{proposition}[Domain-prefix law]
\label{prop:poly-relative-domain-prefix}
For every endoprogram \(U\),
\[
\langle \operatorname{dom}(U)?;U\rangle\varphi
=
\langle U\rangle\varphi.
\]
\end{proposition}

\begin{proof}
\leavevmode
\begin{enumerate}
\item \textbf{Apply the guard-prefix law.}
By the guard-prefix law,
\[
\langle \operatorname{dom}(U)?;U\rangle\varphi
=
\operatorname{dom}(U)\wedge\langle U\rangle\varphi.
\]

\item \textbf{Use monotonicity.}
Since
\[
\langle U\rangle\varphi\le\langle U\rangle\top
=
\operatorname{dom}(U),
\]
the meet is
\[
\langle U\rangle\varphi.
\]
\end{enumerate}
\end{proof}

\subsection{Antidomain predicates}
\label{subsec:poly-relative-antidomain-predicate}

If the predicate fibres are Boolean, define the antidomain predicate by
\[
\operatorname{adom}(U):=\neg\operatorname{dom}(U).
\]

\begin{proposition}[Antidomain annihilation]
\label{prop:poly-relative-antidomain-annihilation}
In Boolean predicate fibres,
\[
\langle \operatorname{adom}(U)?;U\rangle\varphi
=
\bot.
\]
\end{proposition}

\begin{proof}
\leavevmode
\begin{enumerate}
\item \textbf{Apply the guard-prefix law.}
By the guard-prefix law,
\[
\langle \operatorname{adom}(U)?;U\rangle\varphi
=
\operatorname{adom}(U)\wedge\langle U\rangle\varphi.
\]

\item \textbf{Use the domain bound.}
Since
\[
\langle U\rangle\varphi\le\operatorname{dom}(U)
\]
and
\[
\operatorname{adom}(U)=\neg\operatorname{dom}(U),
\]
the meet is bottom.
\end{enumerate}
\end{proof}

\section{Relative guarded Mealy automata}
\label{sec:poly-relative-guarded-automata}

\subsection{Guarded automata}
\label{subsec:poly-relative-guarded-automata-definition}

Let
\[
\Phi:\mathbb A\rightsquigarrow\mathbb B
\]
be a Mealy morphism and let \(Y\) be a right \(\mathbb B\)-action.

Let
\[
(\theta_i)_{i\in I}
\]
be a finite family of state guards on \(S_{\Phi,Y}\),
\[
(R_i)_{i\in I}
\]
a finite family of input-arrow guards on \(A_1\),
\[
(O_i)_{i\in I}
\]
a finite family of output-arrow guards on \(B_1\), and
\[
(T_i)_{i\in I}
\]
a finite family of downstream transition guards on
\[
\left(\int_{\mathbb B}Y\right)_1.
\]

Each quadruple determines a guarded primitive subprogram
\[
G_i:=(\theta_i;R_i/O_i/T_i)^\sharp\hookrightarrow M_{\Phi,Y}.
\]

\begin{definition}[Relative guarded Mealy automaton]
\label{def:poly-relative-guarded-mealy-automaton}
A \textbf{relative guarded Mealy automaton} over \((\Phi,Y)\) is a finite
family of guarded primitive subprograms
\[
G_i=(\theta_i;R_i/O_i/T_i)^\sharp\hookrightarrow M_{\Phi,Y}.
\]
Its one-step program is the finite join
\[
G:=\bigvee_{i\in I}G_i\hookrightarrow M_{\Phi,Y}.
\]
\end{definition}

\subsection{Modal semantics of guarded automata}
\label{subsec:poly-relative-guarded-automata-modal-semantics}

\begin{proposition}[Diamond semantics of guarded automata]
\label{prop:poly-relative-guarded-automaton-diamond}
For a relative guarded Mealy automaton with one-step program
\[
G=\bigvee_{i\in I}G_i,
\]
one has
\[
\langle G\rangle\varphi
=
\bigvee_{i\in I}\langle G_i\rangle\varphi.
\]
\end{proposition}

\begin{proof}
\leavevmode
\begin{enumerate}
\item \textbf{Use the common ambient object.}
All branches are subobjects of the common primitive execution object
\[
M_{\Phi,Y}.
\]

\item \textbf{Apply finite choice.}
The result is the finite choice law for subprograms inside this common ambient
object.
\end{enumerate}
\end{proof}

A relative guarded Mealy automaton is \textbf{witness-disjoint} if
\[
G_i\wedge G_j=\bot
\qquad
(i\ne j).
\]
It is \textbf{branch-disjoint} if
\[
\operatorname{dom}(G_i)\wedge\operatorname{dom}(G_j)=\bot
\qquad
(i\ne j).
\]
It is \textbf{guard-total} if
\[
\bigvee_{i\in I}\operatorname{dom}(G_i)=\top.
\]
When branch-disjointness and guard-totality hold, the branch domains form a
finite disjoint cover of the state object.  The enabled branch is unique at
the level of branch domains, though a branch may still contain multiple
execution witnesses.

\section{Output, downstream, and interface guards}
\label{sec:poly-relative-output-downstream-interface-guards}

\subsection{Output and downstream guards}
\label{subsec:poly-relative-output-downstream-guards}

The relative Mealy setting allows guards to inspect emitted output arrows and
the downstream transitions caused by those outputs.

Let
\[
U\hookrightarrow M_{\Phi,Y}
\]
be a primitive subprogram.

For an output guard
\[
O\hookrightarrow B_1,
\]
define
\[
U_O:=U\wedge\iota_B^\ast O\hookrightarrow M_{\Phi,Y}.
\]

For a downstream transition guard
\[
T\hookrightarrow\left(\int_{\mathbb B}Y\right)_1,
\]
define
\[
U_T:=U\wedge\lambda_{\Phi,Y}^\ast T\hookrightarrow M_{\Phi,Y}.
\]

\subsection{Interface guards}
\label{subsec:poly-relative-interface-guards}

Guards may also be placed directly on response interfaces.  Let
\[
\Gamma\hookrightarrow J_Y=A_0\times Y
\]
be an interface guard.  Pulling it back along
\[
j_{\Phi,Y}:S_{\Phi,Y}\to J_Y
\]
gives a state guard
\[
j_{\Phi,Y}^\ast\Gamma\hookrightarrow S_{\Phi,Y}.
\]
The pullback is displayed by
\[
\begin{tikzcd}[column sep=large,row sep=large]
j_{\Phi,Y}^\ast\Gamma
  \arrow[r]
  \arrow[d,hook]
&
\Gamma
  \arrow[d,hook]
\\
S_{\Phi,Y}
  \arrow[r,"j_{\Phi,Y}"']
&
J_Y
\arrow[ul,phantom,very near start,"\lrcorner"] .
\end{tikzcd}
\]

Thus one may form interface-guarded primitives such as
\[
(j_{\Phi,Y}^\ast\Gamma;R/O/T)^\sharp
\hookrightarrow M_{\Phi,Y}.
\]

\begin{proposition}[External reading of interface guards]
\label{prop:poly-relative-interface-guard-external}
In \(\mathbf{Set}\),
\[
(x,y)\models
\langle j_{\Phi,Y}^\ast\Gamma;R/O/T\rangle\varphi
\]
holds precisely when
\[
\Gamma(p_Px,y)
\]
and there exists an admissible input arrow
\[
a:u\to p_Px
\]
such that
\[
R(a),
\qquad
O(\omega_P(a,x)),
\qquad
T(\omega_P(a,x),y),
\]
and
\[
\left(
\delta_P(a,x),
\sigma(\omega_P(a,x),y)
\right)
\models\varphi.
\]
\end{proposition}

\begin{proof}
\leavevmode
\begin{enumerate}
\item \textbf{Identify the pulled-back state guard.}
The predicate
\[
j_{\Phi,Y}^\ast\Gamma
\]
holds at \((x,y)\) exactly when
\[
\Gamma(p_Px,y)
\]
holds.

\item \textbf{Apply the guarded primitive reading.}
The rest is Proposition~\ref{prop:poly-relative-guarded-primitive-external}
with
\[
\theta=j_{\Phi,Y}^\ast\Gamma.
\]
\end{enumerate}
\end{proof}

\section{The generated guarded program algebra}
\label{sec:poly-relative-generated-program-algebra}

\subsection{Generated programs}
\label{subsec:poly-relative-generated-programs}

\begin{definition}[Relative guarded Mealy program algebra]
\label{def:poly-relative-guarded-program-algebra}
The \textbf{relative guarded Mealy program algebra} generated by
\[
(\Phi,Y)
\]
is the class of endoprograms on
\[
S_{\Phi,Y}
\]
generated from:
\begin{enumerate}
\item tests \(\theta?\) for predicates
\[
\theta\hookrightarrow S_{\Phi,Y};
\]
\item guarded primitive programs
\[
\theta;R/O/T;
\]
\item sequential composition of spans;
\item finite guarded joins of subprograms inside a specified common ambient
span.
\end{enumerate}
\end{definition}

\subsection{Common ambient spans}
\label{subsec:poly-relative-common-ambient-spans}

\begin{remark}[Common ambient spans]
\label{rem:poly-relative-common-ambient-spans}
Finite guarded joins here are joins of subobjects inside a fixed ambient
execution span.  Choice between arbitrary generated spans requires coproducts
of span apexes and is a different, witnessful operation.
\end{remark}

\begin{remark}
\label{rem:poly-relative-program-algebra-from-polynomial}
The polynomial response profile supplies the curried one-step account of
guarded primitives.  Flattening returns the span calculus, where tests,
sequential composition, finite subprogram choice, diamonds, and boxes are
interpreted.
\end{remark}

\section{The one-object specialization}
\label{sec:poly-relative-one-object-specialization}

\subsection{One-object Mealy response profiles}
\label{subsec:poly-relative-one-object-response-profiles}

Let \(M\) and \(N\) be monoid objects in \(\mathcal E\), regarded as
one-object internal categories with the convention
\[
b\circ a=ba.
\]
A Mealy morphism
\[
\Phi:M\rightsquigarrow N
\]
consists of an object \(P\), a right \(M\)-action
\[
x\cdot m:=\delta_P(m,x),
\]
and an \(N\)-valued output cocycle
\[
\omega_P:M\times P\to N
\]
satisfying
\[
\omega_P(1,x)=1,
\qquad
\omega_P(mn,x)=\omega_P(m,x)\omega_P(n,x\cdot m).
\]

Let \(Y\) be a right \(N\)-object.  Then
\[
J_Y=1\times Y\cong Y.
\]
The relative response polynomial has fibre
\[
\mathcal M_{M,N,Y}(X)_y
\cong
\prod_{m\in M}
\sum_{n\in N}
X_{y\cdot n}.
\]
The induced coalgebra is
\[
c_{\Phi,Y}(x,y)(m)
=
\left(
\omega_P(m,x),
\left(
x\cdot m,
y\cdot\omega_P(m,x)
\right)
\right).
\]

\begin{proposition}[One-object cascade response coalgebra]
\label{prop:poly-relative-one-object-cascade-response}
In the one-object case, the relative Mealy response coalgebra is the total
response profile of the cascade action
\[
(x,y)\cdot m
=
\bigl(x\cdot m,\;y\cdot\omega_P(m,x)\bigr).
\]
\end{proposition}

\begin{proof}
\leavevmode
\begin{enumerate}
\item \textbf{Identify the state object.}
Since both object objects are terminal,
\[
S_{\Phi,Y}=P\times Y.
\]

\item \textbf{Identify the response interface.}
The interface is
\[
J_Y\cong Y.
\]

\item \textbf{Specialize the coalgebra.}
The general formula
\[
c_{\Phi,Y}(x,y)(a)
=
\left(
\omega_P(a,x),
\left(
\delta_P(a,x),
\sigma(\omega_P(a,x),y)
\right)
\right)
\]
becomes
\[
c_{\Phi,Y}(x,y)(m)
=
\left(
\omega_P(m,x),
\left(
x\cdot m,
y\cdot\omega_P(m,x)
\right)
\right).
\]
\end{enumerate}
\end{proof}

\subsection{One-object simulations}
\label{subsec:poly-relative-one-object-simulations}

A one-object Mealy simulation
\[
\Theta:\Phi\Rightarrow\Psi
\]
has components
\[
\theta:Q\to P,
\qquad
\alpha:Q\to N.
\]
For every right \(N\)-object \(Y\), the evaluated map is
\[
Q\times Y\to P\times Y,
\qquad
(y,z)\longmapsto(\theta y,z\cdot\alpha_y).
\]
The output-lax comparison component is
\[
\chi(y,z)=\alpha_y.
\]
The simulation equation becomes
\[
\alpha_y\,\omega_P(m,\theta y)
=
\omega_Q(m,y)\,\alpha_{y\cdot m}.
\]

\begin{proposition}[One-object output-lax profile comparison]
\label{prop:poly-relative-one-object-output-lax}
A one-object Mealy simulation evaluates to an output-lax profile comparison
between the corresponding cascade response coalgebras.
\end{proposition}

\begin{proof}
\leavevmode
\begin{enumerate}
\item \textbf{Identify the evaluated state map.}
The state map is
\[
(y,z)\mapsto(\theta y,z\cdot\alpha_y).
\]

\item \textbf{Identify strict input-target preservation.}
The state component of the simulation says
\[
\theta(y\cdot m)=\theta y\cdot m.
\]

\item \textbf{Identify output-lax naturality.}
The output equation
\[
\alpha_y\,\omega_P(m,\theta y)
=
\omega_Q(m,y)\,\alpha_{y\cdot m}
\]
is exactly the one-object output-lax naturality equation.

\item \textbf{Flatten.}
Flattening gives the input-strict functor between the one-object program
categories.
\end{enumerate}
\end{proof}

\section{The stateless specialization}
\label{sec:poly-relative-stateless-specialization}

\subsection{Stateless response coalgebras}
\label{subsec:poly-relative-stateless-response-coalgebras}

Let
\[
F:\mathbb A\to\mathbb B
\]
be an internal functor with object map \(F_0\) and arrow map \(F_1\).  Regard
\(F\) as the representable Mealy morphism with carrier
\[
\begin{tikzcd}
&
A_0
  \arrow[dl,"1_{A_0}"']
  \arrow[dr,"F_0"]
&
\\
A_0 && B_0 .
\end{tikzcd}
\]

For a right \(\mathbb B\)-action \(Y\), the relative state object is
\[
S_{F,Y}:=A_0\times_{F_0,B_0,r}Y.
\]
It is displayed by
\[
\begin{tikzcd}
S_{F,Y}
  \arrow[r,"\pi_Y"]
  \arrow[d,"\pi_A"']
&
Y
  \arrow[d,"r"]
\\
A_0
  \arrow[r,"F_0"']
&
B_0
\arrow[ul,phantom,very near start,"\lrcorner"] .
\end{tikzcd}
\]
The induced coalgebra is
\[
c_{F,Y}:
S_{F,Y}\to
\mathcal M_{\mathbb A,\mathbb B,Y}(S_{F,Y}),
\]
given by
\[
c_{F,Y}(v,y)(a)
=
\left(
F_1(a),
\left(
s_A(a),
\sigma(F_1(a),y)
\right)
\right),
\]
for every arrow \(a:u\to v\).

\begin{proposition}[Stateless relative response coalgebra]
\label{prop:poly-relative-stateless-response-coalgebra}
The coalgebra \(c_{F,Y}\) is the relative response coalgebra associated to the
representable Mealy morphism induced by \(F\).  Its flattened action category
is the action category of the ordinary restricted action \(F^\ast Y\).
\end{proposition}

\begin{proof}
\leavevmode
\begin{enumerate}
\item \textbf{Specialize the carrier.}
For the representable carrier, the state component is the current input object
\(v\in A_0\).

\item \textbf{Identify the state update.}
The state update is forced:
\[
\delta(a,v)=s_A(a).
\]

\item \textbf{Identify the output arrow.}
The emitted output arrow is
\[
\omega(a,v)=F_1(a).
\]

\item \textbf{Substitute into the general coalgebra formula.}
Substituting these into
Definition~\ref{def:poly-relative-mealy-response-coalgebra} gives the
displayed coalgebra.

\item \textbf{Flatten.}
Flattening recovers the usual restriction of a right \(\mathbb B\)-action
along \(F\).
\end{enumerate}
\end{proof}

\subsection{Terminal stateless case}
\label{subsec:poly-relative-terminal-stateless-case}

\begin{corollary}[Terminal stateless coalgebra]
\label{cor:poly-relative-terminal-stateless-coalgebra}
For
\[
Y=\mathbf 1_{\mathbb B},
\]
one has
\[
S_{F,\mathbf 1_{\mathbb B}}\cong A_0,
\]
and the coalgebra becomes
\[
c_F(v)(a)=(F_1(a),s_A(a)).
\]
The recovered program category is
\[
\mathbb A^{\mathrm{op}},
\]
and the output labelling is
\[
F^{\mathrm{op}}:\mathbb A^{\mathrm{op}}\to\mathbb B^{\mathrm{op}}.
\]
\end{corollary}

\begin{proof}
\leavevmode
\begin{enumerate}
\item \textbf{Identify the state object.}
For the terminal downstream action,
\[
S_{F,\mathbf 1_{\mathbb B}}
=
A_0\times_{F_0,B_0,1_{B_0}}B_0
\cong A_0.
\]

\item \textbf{Identify the coalgebra.}
The action of an input arrow \(a:u\to v\) sends \(v\) to \(u=s_A(a)\) and
emits \(F_1(a)\).  Hence
\[
c_F(v)(a)=(F_1(a),s_A(a)).
\]

\item \textbf{Flatten.}
The recovered program category is the action category of the terminal right
\(\mathbb A\)-action, namely
\[
\mathbb A^{\mathrm{op}}.
\]

\item \textbf{Identify the output labelling.}
The output labelling sends \(a\) to \(F_1(a)\), viewed oppositely, hence is
\[
F^{\mathrm{op}}.
\]
\end{enumerate}
\end{proof}

\section{Examples}
\label{sec:poly-relative-examples}

\subsection{Terminal downstream response polynomial}
\label{subsec:poly-relative-example-terminal-downstream-response}

Let
\[
Y=\mathbf 1_{\mathbb B}
=
(B_0\xrightarrow{1_{B_0}}B_0,\tau_B)
\]
be the terminal right \(\mathbb B\)-action.  Then
\[
\tau_B(\beta,t_B\beta)=s_B\beta.
\]
The response interface is
\[
J_{\mathbf 1_{\mathbb B}}=A_0\times B_0.
\]
The fibre formula becomes
\[
\mathcal M_{\mathbb A,\mathbb B,\mathbf 1_{\mathbb B}}(X)_{(v,b)}
\cong
\prod_{a:u\to v}
\sum_{\beta:b'\to b}
X_{(u,b')}.
\]

Moreover,
\[
S_{\Phi,\mathbf 1_{\mathbb B}}
=
P\times_{B_0}B_0
\cong P.
\]
Under this isomorphism,
\[
c_{\Phi,\mathbf 1_{\mathbb B}}(x)(a)
=
(\omega_P(a,x),\delta_P(a,x)).
\]

\begin{corollary}[Terminal Mealy polynomial]
\label{cor:poly-relative-terminal-mealy-polynomial}
The terminal downstream response polynomial of a Mealy morphism
\[
\Phi:\mathbb A\rightsquigarrow\mathbb B
\]
is the special case \(Y=\mathbf 1_{\mathbb B}\) of the relative response
polynomial.  Its induced coalgebra is the terminal output-response profile
\[
x
\longmapsto
\left(
a
\longmapsto(\omega_P(a,x),\delta_P(a,x))
\right).
\]
Flattening this coalgebra recovers the terminal relative program category
\[
\mathsf{Prog}_\Phi.
\]
\end{corollary}

\subsection{Terminal simulation comparison}
\label{subsec:poly-relative-example-terminal-simulation}

\begin{proposition}[Terminal simulation comparison]
\label{prop:poly-relative-terminal-simulation-comparison}
Let
\[
\Theta:\Phi\Rightarrow\Psi
\]
be a Mealy simulation with components \((\theta,\alpha)\).  In the terminal
downstream case, the output-lax profile comparison has state map
\[
\theta:Q\to P
\]
and output-comparison component
\[
\alpha_y:q_P(\theta y)\to q_Q(y).
\]
After flattening, it gives the input-strict functor
\[
\mathsf{Prog}_\Psi\to\mathsf{Prog}_\Phi
\]
and the natural transformation
\[
\Lambda_{\Psi,\mathbf 1_{\mathbb B}}
\Rightarrow
\Lambda_{\Phi,\mathbf 1_{\mathbb B}}H_\Theta
\]
with components
\[
\alpha_y^{\mathrm{op}}:
q_Q(y)\longrightarrow q_P(\theta y)
\]
in \(\mathbb B^{\mathrm{op}}\), where
\[
\alpha_y:q_P(\theta y)\to q_Q(y)
\]
is the underlying arrow of \(\mathbb B\).
\end{proposition}

\begin{proof}
\leavevmode
\begin{enumerate}
\item \textbf{Identify terminal evaluated states.}
In the terminal downstream action,
\[
S_{\Psi,\mathbf 1_{\mathbb B}}\cong Q
\qquad\text{and}\qquad
S_{\Phi,\mathbf 1_{\mathbb B}}\cong P.
\]
The evaluated state map is therefore
\[
\theta:Q\to P.
\]

\item \textbf{Identify the output comparison.}
The output-comparison component remains
\[
\alpha_y:q_P(\theta y)\to q_Q(y)
\]
in \(\mathbb B\).

\item \textbf{View the comparison oppositely.}
Since the terminal downstream action category is
\[
\int_{\mathbb B}\mathbf 1_{\mathbb B}\cong\mathbb B^{\mathrm{op}},
\]
the component is viewed as
\[
\alpha_y^{\mathrm{op}}:
q_Q(y)\to q_P(\theta y)
\]
in \(\mathbb B^{\mathrm{op}}\).

\item \textbf{Flatten.}
Flattening gives the input-strict functor between terminal relative program
categories and the displayed natural transformation.
\end{enumerate}
\end{proof}

\subsection{Relational guarded response semantics}
\label{subsec:poly-relative-example-relational}

In \(\mathbf{Set}\), the response coalgebra
\[
c_{\Phi,Y}(x,y)
\]
assigns to each admissible input arrow \(a:u\to p_Px\) a concrete output arrow
\[
\omega_P(a,x):q_P\delta_P(a,x)\to q_Px
\]
and a concrete continuation state
\[
(\delta_P(a,x),\sigma(\omega_P(a,x),y)).
\]
A guarded existential profile lifting therefore says that some admissible
input position produces an output and continuation satisfying the chosen
guards.  Pulling this lifting back along \(c_{\Phi,Y}\) gives the ordinary
relational diamond of the corresponding restricted primitive execution span.

\subsection{Word inputs and total response profiles}
\label{subsec:poly-relative-example-words}

If \(\mathbb A=I^\ast\) is the one-object category associated to the free
monoid of input words, then the response profile at a state is total over all
words:
\[
w\longmapsto
\left(
\omega_P(w,x),
\delta_P(w,x)
\right).
\]
Letter-step dynamics is obtained either by restricting guarded primitives to
the length-one subpredicate of \(I^\ast\), or by beginning with a generator-level
operation and extending it multiplicatively to \(I^\ast\).

\section{Chapter summary}
\label{sec:poly-relative-summary}

\begin{theorem}[Polynomial response profiles for relative Mealy actions]
\label{thm:poly-relative-main-consolidation}
The constructions of this chapter form the following hierarchy.

\begin{enumerate}
\item A right \(\mathbb A\)-action is an Eilenberg--Moore algebra for the span
monad associated to \(\mathbb A\).

\item When \(\Pi_{t_A}\) exists, a right action admits a response transpose
\[
X\to\mathcal R_{\mathbb A}(X),
\qquad
\mathcal R_{\mathbb A}(X)=\Pi_{t_A}s_A^\ast X,
\]
with fibre
\[
\mathcal R_{\mathbb A}(X)_v
\cong
\prod_{a:u\to v}X_u.
\]

\item For a right \(\mathbb B\)-action \(Y\), the relative response interface is
\[
J_Y=A_0\times Y.
\]

\item The triple \((\mathbb A,\mathbb B,Y)\) determines a relative Mealy
response polynomial
\[
\mathcal M_{\mathbb A,\mathbb B,Y}:
\mathcal E/J_Y\to\mathcal E/J_Y
\]
with fibre
\[
\mathcal M_{\mathbb A,\mathbb B,Y}(X)_{(v,y)}
\cong
\prod_{a:u\to v}
\sum_{\beta:b'\to r(y)}
X_{(u,\sigma(\beta,y))}.
\]

\item Coalgebras
\[
X\to\mathcal M_{\mathbb A,\mathbb B,Y}(X)
\]
are equivalently labelled one-step relative response data.

\item Multiplicative coalgebras are equivalently right \(\mathbb A\)-actions
\(X\) equipped with internal functors
\[
\int_{\mathbb A}X\to\int_{\mathbb B}Y.
\]

\item A Mealy morphism
\[
\Phi:\mathbb A\rightsquigarrow\mathbb B
\]
and a right \(\mathbb B\)-action \(Y\) induce a canonical multiplicative
coalgebra
\[
c_{\Phi,Y}:
S_{\Phi,Y}\to
\mathcal M_{\mathbb A,\mathbb B,Y}(S_{\Phi,Y})
\]
defined by
\[
c_{\Phi,Y}(x,y)(a)
=
\left(
\omega_P(a,x),
\left(
\delta_P(a,x),
\sigma(\omega_P(a,x),y)
\right)
\right).
\]

\item This coalgebra is the dependent-product transpose of the labelled
relative action
\[
(\Phi^\ast Y,\Lambda_{\Phi,Y}).
\]

\item Flattening
\[
c_{\Phi,Y}
\]
recovers the primitive execution span
\[
\mathbf m_{\Phi,Y}
\]
and hence the relative program category
\[
\mathsf{Prog}_{\Phi,Y}
=
\int_{\mathbb A}\Phi^\ast Y.
\]

\item A Mealy simulation
\[
\Theta:\Phi\Rightarrow\Psi
\]
does not generally induce an ordinary coalgebra morphism over
\[
J_Y=A_0\times Y.
\]
Instead, it induces an output-lax profile comparison
\[
c_{\Psi,Y}\longrightarrow c_{\Phi,Y}.
\]

\item Flattening this output-lax profile comparison gives the input-strict,
downstream-lax labelled functor
\[
H_{\Theta,Y}:
\mathsf{Prog}_{\Psi,Y}\to\mathsf{Prog}_{\Phi,Y}
\]
together with the natural transformation
\[
\Lambda_{\Psi,Y}
\Rightarrow
\Lambda_{\Phi,Y}H_{\Theta,Y}.
\]

\item If \(\mathcal E\) is regular, guarded polynomial response liftings pull
back along \(c_{\Phi,Y}\) to guarded span diamonds:
\[
\theta\wedge c_{\Phi,Y}^{\ast}\lambda_{R/O/T}(\varphi)
=
\langle\theta;R/O/T\rangle\varphi.
\]

\item If \(\mathcal E\) is first-order Heyting, the weak universal profile box
pulls back along \(c_{\Phi,Y}\) to the corresponding span box:
\[
c_{\Phi,Y}^{\ast}
\Box^{\mathrm{prof}}_{R/O/T}(\varphi)
=
[R/O/T]\varphi.
\]

\item The strong profile box is generally stricter than the span box, because
it requires every guarded input position to produce a guarded output,
downstream transition, and continuation.

\item Tests, guarded commands, finite guarded choice, domain predicates, and
guarded automata are obtained after flattening on the relative state object
\[
S_{\Phi,Y}.
\]

\item The one-object specialization gives the cascade response coalgebra
\[
c_{\Phi,Y}(x,y)(m)
=
\left(
\omega_P(m,x),
\left(
x\cdot m,
y\cdot\omega_P(m,x)
\right)
\right).
\]

\item The terminal downstream response polynomial is recovered by taking
\[
Y=\mathbf 1_{\mathbb B}.
\]
Its fibre is
\[
\prod_{a:u\to v}
\sum_{\beta:b'\to b}
X_{(u,b')},
\]
and its coalgebra is
\[
x\longmapsto
\left(
a\longmapsto(\omega_P(a,x),\delta_P(a,x))
\right).
\]

\item Representable Mealy morphisms induced by internal functors
\[
F:\mathbb A\to\mathbb B
\]
give stateless relative coalgebras
\[
c_{F,Y}(v,y)(a)
=
\left(
F_1(a),
\left(
s_A(a),
\sigma(F_1(a),y)
\right)
\right).
\]
In the terminal case this reduces to
\[
c_F(v)(a)=(F_1(a),s_A(a)),
\]
and the recovered program category is
\[
\mathbb A^{\mathrm{op}}.
\]
\end{enumerate}
\end{theorem}

\begin{proof}
\leavevmode
\begin{enumerate}
\item \textbf{Actions and response transposes.}
Items (1)--(2) are
Proposition~\ref{prop:poly-relative-actions-as-response-coalgebras}.

\item \textbf{The response polynomial.}
Items (3)--(4) are Definition~\ref{def:poly-relative-response-interface},
Definition~\ref{def:poly-relative-mealy-response-polynomial}, and
Proposition~\ref{prop:poly-relative-fibre-calculation}.

\item \textbf{Coalgebras and labelled actions.}
Item (5) is
Theorem~\ref{thm:poly-relative-currying-labelled-one-step}.  Item (6) is
Proposition~\ref{prop:poly-relative-coalgebras-labelled-actions}.

\item \textbf{Mealy-induced coalgebras and flattening.}
Items (7)--(8) are Definition~\ref{def:poly-relative-mealy-response-coalgebra}
and Theorem~\ref{thm:poly-relative-mealy-coalgebra-transpose}.  Item (9) is
Theorem~\ref{thm:poly-relative-program-category-recovered}.

\item \textbf{Simulation variance.}
Items (10)--(11) are
Theorem~\ref{thm:poly-relative-simulations-profile-comparisons}.

\item \textbf{Profile liftings and boxes.}
Item (12) is
Theorem~\ref{thm:poly-relative-lifting-span-agreement}.  Items (13)--(14) are
Theorem~\ref{thm:poly-relative-profile-box-span-agreement} and
Remark~\ref{rem:poly-relative-weak-strong-profile-boxes}.

\item \textbf{Generated guarded one-step programs.}
Item (15) follows from the flattened span semantics for tests, finite guarded
choice, domain predicates, and guarded automata developed in
Sections~\ref{sec:poly-relative-tests-choice}--\ref{sec:poly-relative-generated-program-algebra}.

\item \textbf{Specializations and examples.}
Item (16) is Proposition~\ref{prop:poly-relative-one-object-cascade-response}.
Item (17) is Corollary~\ref{cor:poly-relative-terminal-mealy-polynomial}.
Item (18) is Proposition~\ref{prop:poly-relative-stateless-response-coalgebra}
and Corollary~\ref{cor:poly-relative-terminal-stateless-coalgebra}.
\end{enumerate}
\end{proof}

\begin{remark}[Final conceptual summary]
\label{rem:poly-relative-final-summary}
The polynomial structure is not imposed externally on the relative Mealy
semantics.  It is obtained by currying the labelled Eilenberg--Moore action
associated to a Mealy morphism and a downstream action:
\[
(\Phi^\ast Y,\Lambda_{\Phi,Y})
\quad\leadsto\quad
c_{\Phi,Y}.
\]
The resulting coalgebra
\[
c_{\Phi,Y}:
S_{\Phi,Y}\to
\mathcal M_{\mathbb A,\mathbb B,Y}(S_{\Phi,Y})
\]
is the total-response presentation of the same one-step data whose uncurried
form is the primitive execution span of
\[
\mathsf{Prog}_{\Phi,Y}.
\]
Mealy simulations appear polynomially as output-lax profile comparisons, not
as ordinary coalgebra morphisms over \(A_0\times Y\).  Flattening sends these
profile comparisons to the labelled functors between relative program
categories.  Pulling guarded polynomial liftings back along the coalgebra
recovers the guarded span modalities.
\end{remark}